\BookUnnumbered{注释}
\begingroup\small\setstretch{1.15}
\NoteChapter{第一章\quad 佟家的奴才}
% 原书页码：311；注释 1:1
\NoteEntry{1}{1}{%
该奏本藏于台湾史语所内阁大库档案，档案编号为038202。内阁大库档案网站有该奏本%
的图像供调阅。%
}
% 原书页码：311；注释 1:2
\NoteEntry{1}{2}{%
关于两位神父在张献忠大西朝的史事，参看：Zürcher, “In the Yellow Tiger's Dcn: Buglio and%
\space%
Magalhāes at the Court of Zhang Xianzhong, 1644-1647", 355-374.%
}
% 原书页码：311；注释 1:3
\NoteEntry{1}{3}{%
关于二位神父到京后的史事，参看：Dunne, Generation of Giants, 327；魏特，《汤若望传》，%
396,401。%
}
% 原书页码：311；注释 1:4
\NoteEntry{1}{4}{%
参看安文思的记录： Magalháes, A New History of China, Containing a Description of the Most%
\space%
Considerable Particulars of That Vast Empire, 346.%
}
% 原书页码：311；注释 1:5
\NoteEntry{1}{5}{%
Oxnam, Ruling from Horseback, 47-49.%
}
% 原书页码：311；注释 1:6
\NoteEntry{1}{6}{%
本文关于两位神父从四川到北京的经历，主要来自两方面材料。一是利类思神父在安文%
思神父1677年去世以后，为他写的一篇简短的传记。在传记开头，利类思自称是“与%
安文思36年未曾分离的挚友，传记于1677年从北京寄出”。后来这篇传记附在安文思%
神父写的《新中国史》书后，于1688年在伦敦出版。第二个材料是，古洛东注解的《圣%
教入川记》。古洛东是法国传教士，原名 Francois Marie Joseph Gourdon。他 19 世纪末年%
人华后在重庆传教，后来在上海从一位耶稣会士那里得到一部手写本，内容是关于利类%
思、安文思明末在成都开辟四川教区的往事。由于古洛东也在四川传教，所以他为这部%
手写本中记录的人物、地点作了注解，最后在1918年出版了该书。该书记录了许多利%
类思、安文思在张献忠朝廷的所见所闻，以及后来他们跟随豪格大军去北京的历史细节。%
Buglio, “An Abridgement of the Life and Death of F. Gabriel Magaillans”, 340-352;古洛东,《圣%
教入川记》。对于两位神父的传教经历，参看：汤开建，《沉与浮：明清鼎革变局中的欧%
洲传教士利类思与安文思》。%
}
% 原书页码：311；注释 1:7
\NoteEntry{1}{7}{%
Buglio, "An Abridgement of the Life and Death of F. Gabriel Magaillans", 345.%
}
% 原书页码：311；注释 1:8
\NoteEntry{1}{8}{%
Buglio, "An Abridgement of the Life and Death of F. Gabriel Magaillans", 346.%
}
% 原书页码：311；注释 1:9
\NoteEntry{1}{9}{%
古洛东，《圣教人川记》，57。%
}
% 原书页码：311；注释 1:10
\NoteEntry{1}{10}{%
古洛东，《圣教人川记》，59。%
}
% 原书页码：311；注释 1:11
\NoteEntry{1}{11}{%
古洛东，《圣教人川记》，52。%
}
% 原书页码：312；注释 1:12
\NoteEntry{1}{12}{%
古洛东，《圣教人川记》，55-56。%
}
% 原书页码：312；注释 1:13
\NoteEntry{1}{13}{%
古洛东，《圣教人川记》，55;Buglio, “An Abridgement of the Life and Death of F. Gabriel%
\space%
Magaillans", 345-346。%
}
% 原书页码：312；注释 1:14
\NoteEntry{1}{14}{%
Crossley, A Translucent Mirror, 140.%
}
% 原书页码：312；注释 1:15
\NoteEntry{1}{15}{%
《满文老档》，165。%
}
% 原书页码：312；注释 1:16
\NoteEntry{1}{16}{%
胡祥雨，《“逃人法”人“顺治律”考——兼谈“逃人法”的应用》。%
}
% 原书页码：312；注释 1:17
\NoteEntry{1}{17}{%
魏特，《汤若望传》，397。%
}
% 原书页码：312；注释 1:18
\NoteEntry{1}{18}{%
Chu, Law and Society in Traditional China, 9, 128, and 280-283; Sommer, Sex, Law, and Society in%
\space%
Late Imperial China, 5-8.%
}
% 原书页码：312；注释 1:19
\NoteEntry{1}{19}{%
对二人转到佟图赖名下，法国学者Pih根据当时传教士寄回罗马的报告，已有详细考证。%
参看： Pih, Le père Gabriel de Magalhães, 94-95.%
}
% 原书页码：312；注释 1:20
\NoteEntry{1}{20}{%
对此问题，20世纪以来，许多清史大家都有论述。认为佟家是汉化的满人的有：孟森，%
《清初三大疑案考实》，120-129；孟森，《明元清系通纪》1:190；陈寅恪，《柳如是别传》，%
1001-1002；牟润孙，《注史斋丛稿》，431-434。认为佟家是汉人，后来迁徙到东北的有：%
郑天挺，《清史探微》，22-28；Crossley，A Translucent Mirror,57-88。%
}
% 原书页码：312；注释 1:21
\NoteEntry{1}{21}{%
陈寅恪，《柳如是别传》，1001。%
}
% 原书页码：312；注释 1:22
\NoteEntry{1}{22}{%
《满文老档》，370；对于清初政治通婚研究，参看：Rawski,“Ch'ing Imperial Marriage and%
\space%
Problems of Rulership".%
}
% 原书页码：312；注释 1:23
\NoteEntry{1}{23}{%
相关内容佟世思在其撰写的家族史中有记载：佟世思，《先高曾祖三世行略》，408-409。%
}
% 原书页码：312；注释 1:24
\NoteEntry{1}{24}{%
Crossley, A Translucent Mirror, 62-63.%
}
% 原书页码：312；注释 1:25
\NoteEntry{1}{25}{%
《清史稿》，2:34。%
}
% 原书页码：312；注释 1:26
\NoteEntry{1}{26}{%
孟森在1936年的文章中就提到，包衣奴才就是身边最近的家奴。包衣奴才只是奴隶的%
一种。相关研究参看：祁美琴，崔灿，《包衣身份再辨》，117-128；孟森，《八旗制度考实》，%
260-262；郑天挺，《清代包衣制度与宦官》；左云鹏，《清代旗下奴仆的地位及其变化》；%
傅克东，《从内佐领和管领谈到清代辛者库人》；杜家骥，《八旗与清朝政治论稿》，435-%
488；鹿智钧，《根本与世仆：清朝旗人的法律地位》，185-194。在红学领域有关家庭内部%
运作的研究中，家中奴仆也是一个部分，这些红学研究对于了解清初奴仆也有贡献。比%
如杜家骥，《清代内务府旗人复杂的旗籍及其多种身份——兼谈曹雪芹家族的旗籍及其%
身份》。另外英文研究中，参看：Elliott,“Vocabulary Notes from the Manchu Archives 2: On%
\space%
the booi"", 18-21; Spence, T’ao Yin; Torbert, The Ch'ing Imperial Household Department; Crossley,%
\space%
A Translucent Mirror, 140-142; Elliott, The Manchu Way, 81–83.%
}
% 原书页码：312；注释 1:27
\NoteEntry{1}{27}{%
魏特，《汤若望传》，468。%
}
% 原书页码：312；注释 1:28
\NoteEntry{1}{28}{%
明代对使用红衣大炮获胜的报告，参看：《明熹宗实录》，68:1a。%
}
% 原书页码：312；注释 1:29
\NoteEntry{1}{29}{%
《明熹宗实录》，70:17b。%
}
% 原书页码：312；注释 1:30
\NoteEntry{1}{30}{%
《明熹宗实录》，79:19b。%
}
% 原书页码：312；注释 1:31
\NoteEntry{1}{31}{%
《皇朝文献通考》，636:461-462。%
}
% 原书页码：312；注释 1:32
\NoteEntry{1}{32}{%
整个行程，参看：Cooper, Rodrigues the Interpreter, 341-342.%
}
% 原书页码：312；注释 1:33
\NoteEntry{1}{33}{%
Cooper, Rodrigues the Interpreter, 338.%
}
% 原书页码：312；注释 1:34
\NoteEntry{1}{34}{%
张小青，《明清之际西洋火炮的输人及其影响》，72。%
}
% 原书页码：313；注释 1:35
\NoteEntry{1}{35}{%
张小青，《明清之际西洋火炮的输入及其影响》，75-76。%
}
% 原书页码：313；注释 1:36
\NoteEntry{1}{36}{%
黄一农，《红夷大炮与皇太极创立的八旗汉军》，77-78。%
}
% 原书页码：313；注释 1:37
\NoteEntry{1}{37}{%
对大凌河之战，魏斐德《洪业：清朝开国史》一书中有精彩论述。%
}
% 原书页码：313；注释 1:38
\NoteEntry{1}{38}{%
《皇朝文献通考》，194:55。%
}
% 原书页码：313；注释 1:39
\NoteEntry{1}{39}{%
《天聪五年八旗值月档（一）》，自《历史档案》，2000（4），4。%
}
% 原书页码：313；注释 1:40
\NoteEntry{1}{40}{%
牟润孙，《注史斋丛稿》，435。%
}
% 原书页码：313；注释 1:41
\NoteEntry{1}{41}{%
对于清朝突然能造出大炮，黄一农认为是满人在1629年突袭中，抓了一批汉人工匠。%
后来清朝铸造的大炮炮身上都刻有这些工匠的名字。同时，黄一农还注意到满人这段时%
间在大海打捞出的沉船中有一尊红衣大炮。于是黄一农认为这批汉人工匠通过模仿海中%
打捞出的大炮，成功仿制出了红衣大炮。我认为黄一农此论对冶炼技术中的技术难点估%
计不足。明末1620年代起，明朝的工匠就在积极仿制，希望制造出红衣大炮来。但是%
炮管等冶炼温度火候掌握不好，造出来的大炮有使用问题及安全隐患。因而明朝在1630%
年代都没有造大炮的技术，需要到澳门购买。所以说，清朝靠自己的能力解决了制造大%
炮中的各种技术问题，这种可能性是很小的。参看：黄一农，《红夷大炮与皇太极创立的%
八旗汉军》。%
}
% 原书页码：313；注释 1:42
\NoteEntry{1}{42}{%
黄一农，《天主教徒孙元化与明末传华的西洋火炮》，911-959。%
}
% 原书页码：313；注释 1:43
\NoteEntry{1}{43}{%
Cooper, Rodrigues the Interpreter, 350.%
}
% 原书页码：313；注释 1:44
\NoteEntry{1}{44}{%
董少新，黄一农，《崇祯年间招募葡兵新考》，78。%
}
% 原书页码：313；注释 1:45
\NoteEntry{1}{45}{%
董少新，黄一农，《崇祯年间招募葡兵新考》，79；黄一农，《红夷大炮与皇太极创立的%
八旗汉军》，90-91。%
}
% 原书页码：313；注释 1:46
\NoteEntry{1}{46}{%
孔有德被封“恭顺王”，参看：《满文老档》，1446。%
}
% 原书页码：313；注释 1:47
\NoteEntry{1}{47}{%
张小青，《明清之际西洋火炮的输入及其影响》，86。%
}
% 原书页码：313；注释 1:48
\NoteEntry{1}{48}{%
《清史稿》，234:9398-9399。%
}
% 原书页码：313；注释 1:49
\NoteEntry{1}{49}{%
关于清朝早期八旗中汉人以及他们成立汉军旗的起源，参看：Li,“The Manchu-Chinese%
\space%
Relationship", 1-32.%
}
% 原书页码：313；注释 1:50
\NoteEntry{1}{50}{%
陈寅恪，《柳如是别传》，158。%
}
% 原书页码：313；注释 1:51
\NoteEntry{1}{51}{%
《皇清开国方略》，341:205。%
}
% 原书页码：313；注释 1:52
\NoteEntry{1}{52}{%
关于地位的变化，参看：Li,“The Rise of the Early Manchu State”, 139-142.%
}
% 原书页码：313；注释 1:53
\NoteEntry{1}{53}{%
关于整个汉军八旗的形成，参看：Li,“The Manchu-Chinese Relationship”, 25.%
}
% 原书页码：313；注释 1:54
\NoteEntry{1}{54}{%
关于红衣大炮数量的估算，参看：黄一农，《红夷大炮与皇太极创立的八旗汉军》，92。%
}
% 原书页码：313；注释 1:55
\NoteEntry{1}{55}{%
《明清史料·乙编》，494。%
}
% 原书页码：313；注释 1:56
\NoteEntry{1}{56}{%
“乌珍超哈”就是汉军的意思。见：《钦定大清会典事例》，1111：1b。%
}
% 原书页码：313；注释 1:57
\NoteEntry{1}{57}{%
Ujen的满语本意，参看：Li,“The Manchu-Chinese Relationship”，25；又陈佳华、傅克东，《八%
旗汉军考略》，19。%
}
% 原书页码：313；注释 1:58
\NoteEntry{1}{58}{%
陈寅恪，《柳如是别传》，158。%
}
% 原书页码：313；注释 1:59
\NoteEntry{1}{59}{%
《大清太宗文皇帝实录》，卷9。%
}
% 原书页码：313；注释 1:60
\NoteEntry{1}{60}{%
《皇清开国方略》，341:262。%
}
% 原书页码：313；注释 1:61
\NoteEntry{1}{61}{%
《清史稿》，列传十八。%
}
% 原书页码：314；注释 1:62
\NoteEntry{1}{62}{%
Dunne, Generation of Giants, 216.%
}
% 原书页码：314；注释 1:63
\NoteEntry{1}{63}{%
Dunne, Generation of Giants, 217。%
}
% 原书页码：314；注释 1:64
\NoteEntry{1}{64}{%
Dunne, Generation of Giants, 318.%
}
% 原书页码：314；注释 1:65
\NoteEntry{1}{65}{%
《火攻挈要》，1643。%
}
% 原书页码：314；注释 1:66
\NoteEntry{1}{66}{%
萧若瑟，《天主教传行中国考》，247。有学者认为此人是尚可喜而非孔有德。%
}
% 原书页码：314；注释 1:67
\NoteEntry{1}{67}{%
孙元化没有听孔有德劝告跟他一起降清。不过，后来孔有德对孙元化一直心存挂念。当%
他得势以后，还主动照顾孙元化的后人。参看：黄一农，《天主教徒孙元化与明末传华的%
西洋火炮》，945。%
}
% 原书页码：314；注释 1:68
\NoteEntry{1}{68}{%
孔有德是孙元化部下，有跟西洋人打交道的经验，最后再救下谢务禄，这里面的联系是%
当时传教士就公认的。参看：《鞑靼战纪》，自《清代西人见闻录》，54。%
}
% 原书页码：314；注释 1:69
\NoteEntry{1}{69}{%
Witek, "Johann Adam Schall von Bell and the Transition from the Ming to the Ch’ing Dynasty",%
\space%
114。%
}
% 原书页码：314；注释 1:70
\NoteEntry{1}{70}{%
满人占领明的领土后，由于八旗中的汉军同时会满汉两种语言，也就成为明疆域内地方%
官差的首选。相关比较参看：Spence,Tsao Yin, 4-5.%
}
% 原书页码：314；注释 1:71
\NoteEntry{1}{71}{%
Witek, "Johann Adam Schall von Bell and the Transition from the Ming to the Ch’ing Dynasty",%
\space%
116。%
}
% 原书页码：314；注释 1:72
\NoteEntry{1}{72}{%
《先高曾祖三世行略》，408-409。%
}
% 原书页码：314；注释 1:73
\NoteEntry{1}{73}{%
《清史稿》，232:9350。%
}
% 原书页码：314；注释 1:74
\NoteEntry{1}{74}{%
Wakeman, The Great Enterprise, 302.%
}
% 原书页码：314；注释 1:75
\NoteEntry{1}{75}{%
礼部大臣对日食观测结果的报告，见：《西洋新法历书·奏疏》，1258；《正教奉褒》，24b;%
\space%
Witek, "Johann Adam Schall von Bell and the Transition from the Ming to the Ch’ing Dynasty", 114.%
}
% 原书页码：314；注释 1:76
\NoteEntry{1}{76}{%
Witek, "Johann Adam Schall von Bell and the Transition from the Ming to the Ch’ing Dynasty",%
\space%
116.%
}
% 原书页码：314；注释 1:77
\NoteEntry{1}{77}{%
魏特，《汤若望传》，244。%
}
% 原书页码：314；注释 1:78
\NoteEntry{1}{78}{%
陈寅恪，《柳如是别传》，1001。%
}
% 原书页码：314；注释 1:79
\NoteEntry{1}{79}{%
Hummel, Eminent Chinese of the Ching Period (1644–1912), 793; Crossley, A Translucent Mirror, 110.%
}
% 原书页码：314；注释 1:80
\NoteEntry{1}{80}{%
方豪；《中国天主教史人物传》2:50-52。%
}
% 原书页码：314；注释 1:81
\NoteEntry{1}{81}{%
黄一农，《天主教徒孙元化与明末传华的西洋火炮》，947。%
}
% 原书页码：314；注释 1:82
\NoteEntry{1}{82}{%
Hummel, Eminent Chinese of the Ch’ing Period (1644–1912), 793.%
}
% 原书页码：314；注释 1:83
\NoteEntry{1}{83}{%
Pih, Le père Gabriel de Magalháes, 95-96；魏特，《汤若望传》，468-469。%
}
% 原书页码：314；注释 1:84
\NoteEntry{1}{84}{%
《耶稣会士中国书简集》3:26。%
}
% 原书页码：314；注释 1:85
\NoteEntry{1}{85}{%
魏特，《汤若望传》，468-469。具体关于二位神父去除奴籍以后，变成了哪一种身份，%
还有待确认。根据现知清初旧例，出了奴籍以后的人被称为“开户”。关于这类人，参看：%
刘小萌，《八旗户籍中的旗下人诸名称考释》，65-66。%
}
% 原书页码：314；注释 1:86
\NoteEntry{1}{86}{%
佟家引荐了两位神父给顺治皇帝。安文思神父为顺治皇帝管理他私人所有的西洋奇珍异%
物。Buglio, "An Abridgement of the Life and Death of F. Gabriel Magaillans”, 346.%
}
% 原书页码：314；注释 1:87
\NoteEntry{1}{87}{%
根据满文档案，顺治皇帝派佟吉负责教堂修建的事务。参看：《清初西洋传教士满文档案%
译本》，266。%
}
% 原书页码：315；注释 1:88
\NoteEntry{1}{88}{%
也不是所有的清朝奴仆都想去除奴籍，这很大程度上取决于其主人的经济条件以及%
社会地位等诸多因素。一方面，有很多奴仆由于生活艰难，私自逃跑；但同时也有一%
些汉人的平民，发现给满人贵族为奴的生活更富裕，主动加人奴籍。对于这个问题%
的探讨，参看：杜家骥，《八旗与清朝政治论稿》，435-488；鹿智钧，《根本与世仆》，%
185-194。%
}
% 原书页码：315；注释 1:89
\NoteEntry{1}{89}{%
比如曹寅一家，从他爷爷起，曹家已经去除奴籍，但曹家知道，他们家世中最显赫的部%
分就是做过爱新觉罗家的家奴。所以到曹寅一代，他一直在康熙面前称奴，而他也确实%
因为有康熙家包衣奴才的身份，取得了很多特权和利益。参看：傅克东，《八旗户籍制度%
初探》, 34-35; Spence, Tsao Yin, 10-11。%
}
% 原书页码：315；注释 1:90
\NoteEntry{1}{90}{%
Dunne, Generation of Giants, 326.%
}
% 原书页码：315；注释 1:91
\NoteEntry{1}{91}{%
波士顿大学梅欧金教授为我提供了该信的原件以及英文翻译。该信的原始出处：%
Longobardi, Niccolò SJ. "Originale [in Portuguese] della lettera del P. Nicolao Longobardi scritta%
\space%
da Pechino al P. Francesco Brancato consultor della Viceprovincia della China". Ms., Beijing, 17%
\space%
August 1651. ARSI, Japonica Sinica vol. 142, f. 135r-136v. 这封信的原件是葡萄牙文写成%
的，当时还有一份意大利文的翻译，出处是： Longobardi, Niccolò SJ.“Copia [in Italian]%
della lettera del P. Nicolao Longobardi scritta da Pechino al P. Francesco Brancato consultor della%
\space%
Viceprovincia della China". Ms., Beijing, 17 August 1651. ARSI, Japonica Sinica vol. 142, ff.%
\space%
137r-138v.%
}
% 原书页码：315；注释 1:92
\NoteEntry{1}{92}{%
Pih, Le père Gabriel de Magalhães, chapter 3. 关于双方恩怨，参看：夏伯嘉，《从安文思与汤%
若望之间的恩怨论说明宫廷耶稣会士》。%
}
% 原书页码：315；注释 1:93
\NoteEntry{1}{93}{%
Dunne, Generation of Giants, 325-338.%
}
% 原书页码：315；注释 1:94
\NoteEntry{1}{94}{%
Buglio, "An Abridgement of the Life and Death of F. Gabriel Magaillans", 345.%
}
% 原书页码：315；注释 1:95
\NoteEntry{1}{95}{%
古洛东，《圣教人川记》，52。%
}
% 原书页码：315；注释 1:96
\NoteEntry{1}{96}{%
《皇清开国方略》，341:262。%
}
% 原书页码：315；注释 1:97
\NoteEntry{1}{97}{%
Givon, "Connecting Eurasia: Jesuit Experimentation with Overland Mobility Between China and%
\space%
Europe, 1656–1664”, 655–656.%
}
\NoteChapter{第二章\quad 满汉之争中的传教士}
% 原书页码：315；注释 2:1
\NoteEntry{2}{1}{%
曹树基，《中国人口史（第五卷清时期）》，17。%
}
% 原书页码：315；注释 2:2
\NoteEntry{2}{2}{%
Brockey, Journey to the East, 112.%
}
% 原书页码：315；注释 2:3
\NoteEntry{2}{3}{%
信件的英文译本，参看： Chan, “A European Document on the Fall of the Ming Dynasty (1644-%
1649)", 105.%
}
% 原书页码：315；注释 2:4
\NoteEntry{2}{4}{%
Brockey, Journey to the East, 112.%
}
% 原书页码：315；注释 2:5
\NoteEntry{2}{5}{%
Rouleau, "The First Chinese Priest of the Society of Jesus", 3-50.%
}
% 原书页码：315–316；注释 2:6
\NoteEntry{2}{6}{%
关于汤若望与顺治的私交早有研究。比较著名的是陈垣的《汤若望与木陈忞》。该文比%
较了汤若望、禅师木陈忞与顺治之间的不同私交。传教士安文思与顺治的私人关系少有%
人注意。当时世界的航海贸易已经成熟，西洋的各种精巧器件、把玩之物已经大量进人%
中国权贵家中。顺治作为皇帝，当然也有各种西洋送来的奇珍异宝。安文思神父是个能%
工巧匠，能够修理制作很多东西。所以他被安排管理和维修顺治的这些宝贝。从跟顺治%
的关系上来说，算是顺治私下的仆人。安文思去世以后，利类思神父为其写的小传中曾%
专门提到他的这个角色。另外，康熙为安文思写的墓志铭上也专门提到了安文思为他父%
亲效力的事。安文思的墓碑现存于中共北京市委党校（北京行政学院）内滕公栅栏传教%
士墓地，该墓地现为全国重点文物保护单位。%
}
% 原书页码：316；注释 2:7
\NoteEntry{2}{7}{%
关于卫匡国神父在福建的情况简介，参看：Menegon, Ancestors, Virgins, and Friars, 94-95.%
}
% 原书页码：316；注释 2:8
\NoteEntry{2}{8}{%
陈垣，《汤若望与木陈忞》。%
}
% 原书页码：316；注释 2:9
\NoteEntry{2}{9}{%
魏特，《汤若望传》，323-324。%
}
% 原书页码：316；注释 2:10
\NoteEntry{2}{10}{%
顺治出家，以及五台山的相关传说考证，参看：孟森，《清初三大疑案考实》，12-57。%
}
% 原书页码：316；注释 2:11
\NoteEntry{2}{11}{%
谢和耐《中国和基督教》应该算是从中西文化冲突角度观察明清基督教的奠基之作。%
在谢作发表的1980年以前，历狱被看作一桩在中国发生的排外案子。相关著作参阅：%
Cohen, China and Christianity, 20-34; Young, "An Early Confucian Attack on Christianity", 159–%
186; Young, Confucianism and Christianity, 77–96.%
\space%
1980年代后，整个历狱案子主要在中西文化冲突的框架下进行解读，相关论文参看：%
Zurndorfer, "One Adam Having Driven Us out of Paradise, Another Has Driven Us out of China”,%
\space%
141-168; Chu, "Scientific Dispute in the Imperial Court: The 1664 Calendar Case", 7-34.%
}
% 原书页码：316；注释 2:12
\NoteEntry{2}{12}{%
1990年代以来，黄一农教授发表了多篇论文，深人研究了历狱涉及的多个方面，发现%
了许多前人未发现的细节。黄一农认为案子涉及很多实际问题，并非是在空谈文化冲%
突。但是黄一农教授并未探索满人在这个案子中扮演的角色。对于以往历狱研究的综%
述，参看 Jami,“Revisiting the Calendar Case(1664-1669)”。黄一农，《择日之争与“康熙%
历狱”》,247-280;《清前期对觜、参两宿先后次序的争执：社会天文学史之一个案研究》，%
71-93。%
}
% 原书页码：316；注释 2:13
\NoteEntry{2}{13}{%
关于议政王大臣会议在清初的地位和作用，参看：傅宗懋，《清初议政体制之研究》，%
245-294。另英文研究参看：Oxnam,“Policies and Institutions of the Oboi Regency, 1661-%
1669", 271-273.%
}
% 原书页码：316；注释 2:14
\NoteEntry{2}{14}{%
《清史稿》，卷5。%
}
% 原书页码：316；注释 2:15
\NoteEntry{2}{15}{%
Spence, "The Kang-hsi Reign", 131.%
}
% 原书页码：316；注释 2:16
\NoteEntry{2}{16}{%
杨光先1659年的参劾，见黄伯禄，《正教奉褒》，33b。%
}
% 原书页码：316；注释 2:17
\NoteEntry{2}{17}{%
奥地利国家博物馆现存1661年版。另，该书的各种版本比较，参看：纪建勋，《〈碑记赠%
言合刻〉版本篇目互校》，81-97。%
}
% 原书页码：316；注释 2:18
\NoteEntry{2}{18}{%
1663年开封犹太教教堂外立的石碑，就是一位汉人进士撰写的。碑文相关研究参看：%
Leslie, The Survival of the Chinese Jews, 40-41.%
}
% 原书页码：316；注释 2:19
\NoteEntry{2}{19}{%
前面提到的龚鼎孳就把汤若望与僧一行作对比。南怀仁也把汤研治天文与唐代的僧一行%
做类比。参看：《熙朝定案》，1394；黄伯禄，《正教奉褒》，54b。%
}
% 原书页码：316；注释 2:20
\NoteEntry{2}{20}{%
魏特，《汤若望传》，323-324。%
}
% 原书页码：316；注释 2:21
\NoteEntry{2}{21}{%
陆世仪，《复社纪略》，131。%
}
% 原书页码：316；注释 2:22
\NoteEntry{2}{22}{%
王士禛，《池北偶谈》，88。%
}
% 原书页码：316；注释 2:23
\NoteEntry{2}{23}{%
黄仁宇《万历十五年》的第二、三章中对于明末朝廷中的参劾文化有细致描述。另外，%
关于明代的监察制度，参看：Hucker, The Censorial System of Ming China.%
}
% 原书页码：316；注释 2:24
\NoteEntry{2}{24}{%
王士禛，《池北偶谈》，88。%
}
% 原书页码：317；注释 2:25
\NoteEntry{2}{25}{%
黄一农，《张宸生平及其与杨光先间的冲突》，82-83。%
}
% 原书页码：317；注释 2:26
\NoteEntry{2}{26}{%
《四库全书总目提要》，卷133。%
}
% 原书页码：317；注释 2:27
\NoteEntry{2}{27}{%
《圣祖仁皇帝实录》，卷15。%
}
% 原书页码：317；注释 2:28
\NoteEntry{2}{28}{%
黄一农，《清前期对觜、参两宿先后次序的争执：社会天文学史之一个案研究》，71-94。%
审讯汤若望时，有问到这个问题。见《清初西洋传教士满文档案译本》，164。%
}
% 原书页码：317；注释 2:29
\NoteEntry{2}{29}{%
黄一农，《张宸生平及其与杨光先间的冲突》，82-83。%
}
% 原书页码：317；注释 2:30
\NoteEntry{2}{30}{%
朱维铮，《走出中世纪》，205。%
}
% 原书页码：317；注释 2:31
\NoteEntry{2}{31}{%
对于议政王大臣会议的人员组成和背景分析，参看：Oxnam,“Policies and Institutions of the%
\space%
Oboi Regency, 1661–1669", 271-273.%
}
% 原书页码：317；注释 2:32
\NoteEntry{2}{32}{%
魏特，《汤若望传》，501。%
}
% 原书页码：317；注释 2:33
\NoteEntry{2}{33}{%
《西洋新法历书·奏疏》，1300。%
}
% 原书页码：317；注释 2:34
\NoteEntry{2}{34}{%
汤若望在中国做占算并不是什么秘密，当时许多在华传教士都不同意汤若望参与这些占%
算活动。有传教士把汤若望的这些活动报告给了罗马教廷，希望罗马方面能出面制止。%
参看： Collani, “Astronomy versus Astrology", 439.%
回欧洲汇报中国事务的卫匡国神父为汤若望辩护，认为汤若望参与的占算和欧洲当%
时流行的占星术是一样的。既然罗马允许欧洲的占星术，那么汤若望在中国参与占算也%
应该被允许。更为关键的是，卫匡国神父认为汤若望的占算活动总体上是有利于天主教%
在华传播的。Collani,“Two Astronomers”,83-84。%
}
% 原书页码：317；注释 2:35
\NoteEntry{2}{35}{%
《明清史料·甲编》，317。这段史料西方学者很早就注意到了，并探索过到底汤若望奏%
疏中影射的是哪个女人。见：Rosso, Apostolic Legations to China of the Eighteenth Century, 116-%
118; Lippiello, "Astronomy and Astrology: Johann Adam Schall von Bell", 424-425.%
}
% 原书页码：317；注释 2:36
\NoteEntry{2}{36}{%
汤若望掌控的钦天监总共上报了22次日珥。与之相对，后来的南怀仁神父领导的钦天%
监一次日珥都没有上报过。《世祖章皇帝实录》，卷5，卷12，卷21，卷30，卷35，卷%
70，卷72，卷75，卷76，卷78，卷79，卷80，卷84，卷86，卷91，卷94，卷98，%
卷100，卷118，卷129，卷132。%
}
% 原书页码：317；注释 2:37
\NoteEntry{2}{37}{%
魏特，《汤若望传》，396，455。%
}
% 原书页码：317；注释 2:38
\NoteEntry{2}{38}{%
Collani, "Astronomy versus Astrology", 439.%
}
% 原书页码：317；注释 2:39
\NoteEntry{2}{39}{%
五世达赖的自传中记载汤若望，纯粹是惊讶于他的预测技术，他和汤若望没有私交，对%
他也没有其他评价。五世达赖喇嘛，《五世达赖喇嘛传：云裳》，328。%
}
% 原书页码：317；注释 2:40
\NoteEntry{2}{40}{%
《世祖章皇帝实录》，卷68。《实录》中记录这件事时，就是以满汉官员互相争执的形式。%
}
% 原书页码：317；注释 2:41
\NoteEntry{2}{41}{%
《世祖章皇帝实录》，卷69。%
}
% 原书页码：317；注释 2:42
\NoteEntry{2}{42}{%
最早全面梳理历狱案子的是魏特的《汤若望传》。在书中，魏特认为汤若望没有参与到%
测算中，他只是在做科学的天文观测。%
}
% 原书页码：317；注释 2:43
\NoteEntry{2}{43}{%
Deiwiks, "The Secret Manchu Documents on the Trial of Jesuit Missionary Johann Adam Schall%
(1592-1666) before the Supreme Court of Peking", 643.%
}
% 原书页码：317；注释 2:44
\NoteEntry{2}{44}{%
黄一农，《择日之争与“康熙历狱”》，257-258。%
}
% 原书页码：317–318；注释 2:45
\NoteEntry{2}{45}{%
历狱的审判资料是用满文记录的，现存中国第一历史档案馆。2015年全部档案已经翻%
译成汉语出版。参看《清初西洋传教士满文档案译本》。对该档案的研究有：安双成，%
《汤若望案始末》;Deiwiks,“The Secret Manchu Documents";Deiwiks,“Some Cultural and%
\space%
Psychological Aspects of the Trial of Johann Adam Schall Before the Supreme Court of Peking%
\space%
According to the Secret Manchu Documents".%
}
% 原书页码：318；注释 2:46
\NoteEntry{2}{46}{%
《清初西洋传教士满文档案译本》，44。%
}
% 原书页码：318；注释 2:47
\NoteEntry{2}{47}{%
《清初西洋传教士满文档案译本》，66。%
}
% 原书页码：318；注释 2:48
\NoteEntry{2}{48}{%
《清初西洋传教士满文档案译本》，68。汤若望在他执掌钦天监的二十年中，被参劾过很%
多次。但他以前都是引经据典，以专家自居。这是第一次他以“没有参与”“不懂”作%
为理由来抗辩。对于汤若望其他参劾抗辩，参看：《西洋新法历书·奏疏》，1314-1319。%
}
% 原书页码：318；注释 2:49
\NoteEntry{2}{49}{%
《清初西洋传教士满文档案译本》，189。%
}
% 原书页码：318；注释 2:50
\NoteEntry{2}{50}{%
《清初西洋传教士满文档案译本》，189。%
}
% 原书页码：318；注释 2:51
\NoteEntry{2}{51}{%
《清初西洋传教士满文档案译本》，184。%
}
% 原书页码：318；注释 2:52
\NoteEntry{2}{52}{%
《清初西洋传教士满文档案译本》，217。%
}
% 原书页码：318；注释 2:53
\NoteEntry{2}{53}{%
《清初西洋传教士满文档案译本》，179-180。%
}
% 原书页码：318；注释 2:54
\NoteEntry{2}{54}{%
《清初西洋传教士满文档案译本》，267-268。%
}
% 原书页码：318；注释 2:55
\NoteEntry{2}{55}{%
《清初西洋传教士满文档案译本》，266。%
}
% 原书页码：318；注释 2:56
\NoteEntry{2}{56}{%
叶梦珠，《阅世编》，10。%
}
% 原书页码：318；注释 2:57
\NoteEntry{2}{57}{%
其实传教士早就预见到朝廷政治斗争的风险。在1649年，由安文思起草、十位传教士%
签名的一封信就警告过汤若望，提醒他在朝廷做官，管理一个部门，不可避免会得罪人，%
早晚会有政敌出现。魏特，《汤若望传》，477。%
}
% 原书页码：318；注释 2:58
\NoteEntry{2}{58}{%
魏特，《汤若望传》，516。%
}
% 原书页码：318；注释 2:59
\NoteEntry{2}{59}{%
Gernet, China and the Christian Impact, 43.%
}
% 原书页码：318；注释 2:60
\NoteEntry{2}{60}{%
Jami, The Emperor's New Mathematics, 58.%
}
\NoteChapter{第三章\quad 从铲除鳌拜到《尼布楚条约》谈判}
% 原书页码：318；注释 3:1
\NoteEntry{3}{1}{%
昭梿，《啸亭杂录》，476；孟昭信，《康熙大帝全传》，27-28。%
}
% 原书页码：318；注释 3:2
\NoteEntry{3}{2}{%
杨珍，《康熙朝鳌拜罪案辨析》，89-92。%
}
% 原书页码：318；注释 3:3
\NoteEntry{3}{3}{%
《满文老档》，697。%
}
% 原书页码：318；注释 3:4
\NoteEntry{3}{4}{%
Rawski, "Ch'ing Imperial Marriage and Problems of Rulership”, 177-179.%
}
% 原书页码：318；注释 3:5
\NoteEntry{3}{5}{%
杜家骥总结了清初所有的满蒙联姻。参看：杜家骥，《清朝满蒙联姻研究》，596-624。%
}
% 原书页码：318；注释 3:6
\NoteEntry{3}{6}{%
1644年人关以前，满蒙通婚85次，其中33次也就是近40\%都是和科尔沁部展开的。%
杜家骥，《清朝满蒙联姻研究》，12。%
}
% 原书页码：318；注释 3:7
\NoteEntry{3}{7}{%
顺治后宫研究，参看：孟森，《清世祖董鄂妃生死特殊典礼》，178-184；《世祖出家事考%
实》，237-247。%
}
% 原书页码：318；注释 3:8
\NoteEntry{3}{8}{%
蒋兆成，王日根，《康熙传》，9。%
}
% 原书页码：319；注释 3:9
\NoteEntry{3}{9}{%
杨珍，《苏麻喇姑研究》。%
}
% 原书页码：319；注释 3:10
\NoteEntry{3}{10}{%
Spence, "The Kang-hsi Reign", 140-141.%
}
% 原书页码：319；注释 3:11
\NoteEntry{3}{11}{%
Oxnam, Ruling from Horseback, 169.%
}
% 原书页码：319；注释 3:12
\NoteEntry{3}{12}{%
Hummel, Eminent Chinese of the Ch'ing Period (1644–1912), 663.%
}
% 原书页码：319；注释 3:13
\NoteEntry{3}{13}{%
对于中国传统士人而言，“文正”是谥号中的最高荣誉。而对于满人来说，“忠”才是他%
们最看重的，所以“文忠”是他们认为的最高评价。清初四朝，“文正”给出了四次。%
而“文忠”只给出一次。后来清末慈禧颁给李鸿章的谥号是“文忠”，其实是对李作为%
汉臣的最高认可。参看：宋秉仁，《从文臣谥号看明清谥法理论》，255。%
}
% 原书页码：319；注释 3:14
\NoteEntry{3}{14}{%
Lawrence Kessler 注意到了相关的上奏，并做了统计。参看：Kessler, K'ang-hsi and the Con-%
solidation of Ching Rule, 1661–1684, 188.%
}
% 原书页码：319；注释 3:15
\NoteEntry{3}{15}{%
Kessler, K'ang-hsi and the Consolidation of Ch’ing Rule, 1661–1684, 57.%
}
% 原书页码：319；注释 3:16
\NoteEntry{3}{16}{%
常江，李理，《清宫大内侍卫》。%
}
% 原书页码：319；注释 3:17
\NoteEntry{3}{17}{%
常江，李理，《清宫大内侍卫》，134。%
}
% 原书页码：319；注释 3:18
\NoteEntry{3}{18}{%
常江，李理，《清宫大内侍卫》，149-150。%
}
% 原书页码：319；注释 3:19
\NoteEntry{3}{19}{%
Hummel, Eminent Chinese of the Ch’ing Period (1644–1912), 795.%
}
% 原书页码：319；注释 3:20
\NoteEntry{3}{20}{%
《圣祖仁皇帝圣训》，411:169。%
}
% 原书页码：319；注释 3:21
\NoteEntry{3}{21}{%
《北京图书馆藏中国历代石刻拓本汇编》，65:107。%
}
% 原书页码：319；注释 3:22
\NoteEntry{3}{22}{%
南怀仁参劾的折子全文，参看：《熙朝定案》，1389-1390。%
}
% 原书页码：319；注释 3:23
\NoteEntry{3}{23}{%
《圣祖仁皇帝实录》，卷27。%
}
% 原书页码：319；注释 3:24
\NoteEntry{3}{24}{%
《圣祖仁皇帝实录》，卷27。%
}
% 原书页码：319；注释 3:25
\NoteEntry{3}{25}{%
Kessler, K ang-hsi and the Consolidation of Ch’ing Rule, 1661–1684, 63.%
}
% 原书页码：319；注释 3:26
\NoteEntry{3}{26}{%
Verbiest, The Astronomia Europaea of Ferdinand Verbiest, S.J, 66.%
}
% 原书页码：319；注释 3:27
\NoteEntry{3}{27}{%
Cullen, C., \& Jami, C. "Christmas 1668 and After: How Jesuit Astronomy Was Restored to Power%
\space%
in Beijing", 10-11; Pih, Le père Gabriel de Magalhāes, 214–215.%
}
% 原书页码：319；注释 3:28
\NoteEntry{3}{28}{%
《熙朝定案》，1390。%
}
% 原书页码：319；注释 3:29
\NoteEntry{3}{29}{%
《正教奉褒》，58。%
}
% 原书页码：319；注释 3:30
\NoteEntry{3}{30}{%
《圣祖仁皇帝实录》，卷30。%
}
% 原书页码：319；注释 3:31
\NoteEntry{3}{31}{%
《圣祖仁皇帝实录》，卷30。%
}
% 原书页码：319；注释 3:32
\NoteEntry{3}{32}{%
汤若望墓现存于滕公栅栏，位于车公庄大街6号中共北京市委党校（北京行政学院）院%
内。很多著名来华传教士都安葬于此。这里是全国重点文物保护单位。%
}
% 原书页码：319；注释 3:33
\NoteEntry{3}{33}{%
Verbiest, The Astronomia Europaea of Ferdinand Verbiest, S. J., 61.%
}
% 原书页码：319；注释 3:34
\NoteEntry{3}{34}{%
Verbiest, The Astronomia Europaea of Ferdinand Verbiest, S. J., 66.%
}
% 原书页码：319；注释 3:35
\NoteEntry{3}{35}{%
Verbiest, The Astronomia Europaea of Ferdinand Verbiest, S. J., 66, 68-69.%
}
% 原书页码：319；注释 3:36
\NoteEntry{3}{36}{%
Verbiest, The Astronomia Europaea of Ferdinand Verbiest, S. J.%
}
% 原书页码：319；注释 3:37
\NoteEntry{3}{37}{%
Verbiest, The Astronomia Europaea of Ferdinand Verbiest, S. J., 63.%
}
% 原书页码：319；注释 3:38
\NoteEntry{3}{38}{%
《正教奉褒》，57a-58b。%
}
% 原书页码：320；注释 3:39
\NoteEntry{3}{39}{%
《清初西洋传教士满文档案译本》，179-180；另见魏特，《汤若望传》，493。%
}
% 原书页码：320；注释 3:40
\NoteEntry{3}{40}{%
当时谋反主谋已经供出钱谦益。柳如是贿赂了南方相关官员后，官员上书帮助钱谦益脱%
罪。脱罪的理由是钱谦益不认识主谋。陈寅恪考证出钱谦益跟案件确实有牵连。陈寅恪，%
《柳如是别传》，913-923。%
}
% 原书页码：320；注释 3:41
\NoteEntry{3}{41}{%
《正教奉褒》，65a-65b。%
}
% 原书页码：320；注释 3:42
\NoteEntry{3}{42}{%
《正教奉褒》，33b-34a。%
}
% 原书页码：320；注释 3:43
\NoteEntry{3}{43}{%
魏特，《汤若望传》，468。%
}
% 原书页码：320；注释 3:44
\NoteEntry{3}{44}{%
陈垣，《汤若望与木陈忞》，502。%
}
% 原书页码：320；注释 3:45
\NoteEntry{3}{45}{%
《正教奉褒》，72a。%
}
% 原书页码：320；注释 3:46
\NoteEntry{3}{46}{%
Verbiest, The Astronomia Europaea of Ferdinand Verbiest, S. J., 99.%
}
% 原书页码：320；注释 3:47
\NoteEntry{3}{47}{%
Verbiest, The Astronomia Europaea of Ferdinand Verbiest, S. J., 101.%
}
% 原书页码：320；注释 3:48
\NoteEntry{3}{48}{%
《熙朝定案》，1394。%
}
% 原书页码：320；注释 3:49
\NoteEntry{3}{49}{%
《熙朝定案》，1395。%
}
% 原书页码：320；注释 3:50
\NoteEntry{3}{50}{%
Landry-Deron, "The Kangxi Emperor's Lessons in Western Sciences as Recounted by the Jesuit%
\space%
Fathers J. Bouvet and J.-F. Gerbillon", 257-270.%
}
% 原书页码：320；注释 3:51
\NoteEntry{3}{51}{%
《圣祖仁皇帝实录》，卷146。%
}
% 原书页码：320；注释 3:52
\NoteEntry{3}{52}{%
Verbiest, "Journeys into Tartary", 116.%
}
% 原书页码：320；注释 3:53
\NoteEntry{3}{53}{%
舒理广，胡建中，周铮，《南怀仁与中国清代铸造的大炮》，25-31。%
}
% 原书页码：320；注释 3:54
\NoteEntry{3}{54}{%
Stary, "The Manchu Cannons' Cast by Ferdinand Verbiest and the Hitherto Unknown Title of%
\space%
His Instructions", 215-225.%
}
% 原书页码：320；注释 3:55
\NoteEntry{3}{55}{%
Verbiest, Astronomia Europaea sub imperatore Tartaro Sinico Cám Hý appellato ex umbra in lucem%
\space%
revocata.%
}
% 原书页码：320；注释 3:56
\NoteEntry{3}{56}{%
康熙的大哥应该是福全，他和康熙以及佟家都保持着良好关系。Verbiest, Astronomia%
\space%
Europaea sub imperatore Tartaro Sinico Cám Hý appellato ex umbra in lucem revocata, 87–88.%
}
% 原书页码：320；注释 3:57
\NoteEntry{3}{57}{%
Scheel, "Beijing Precursor”, 245-270; Needham, Science and Civilization in China, 4: 225.%
}
% 原书页码：320；注释 3:58
\NoteEntry{3}{58}{%
关于欧洲当时的中国热，参看：Hsia, Chinesia: The European Construction of China in the%
\space%
Literature of the 17th and 18th Centuries.%
}
% 原书页码：320；注释 3:59
\NoteEntry{3}{59}{%
Wills, Embassies and Illusions, 120.%
}
% 原书页码：320；注释 3:60
\NoteEntry{3}{60}{%
《鞑靼旅行记》，79。%
}
% 原书页码：320；注释 3:61
\NoteEntry{3}{61}{%
《鞑靼旅行记》，76。%
}
% 原书页码：320；注释 3:62
\NoteEntry{3}{62}{%
Hao, "Ferdinand Verbiest and His Role in the Formation of Sino-Russian Relations”, 284-285.%
}
% 原书页码：320；注释 3:63
\NoteEntry{3}{63}{%
Loewenthal, "Nikolai Gavrilovich Spafarii-Milesku (1636-1708): A Biobibliography", 95-111.%
}
% 原书页码：320；注释 3:64
\NoteEntry{3}{64}{%
Mancall, Russia and China, 82–94.%
}
% 原书页码：320；注释 3:65
\NoteEntry{3}{65}{%
Huang, 1587, The Year of No Significance: the Ming Dynasty in Decline, 85, 92.%
}
% 原书页码：320；注释 3:66
\NoteEntry{3}{66}{%
Givon, “Connecting Eurasia: Jesuit Experimentation with Overland Mobility Between China and%
\space%
Europe, 1656-1664".%
}
% 原书页码：321；注释 3:67
\NoteEntry{3}{67}{%
Sebes, The Jesuits and the Sino-Russian Treaty of Nerchinsk, 96.%
}
% 原书页码：321；注释 3:68
\NoteEntry{3}{68}{%
Szczesniak, "Diplomatic Relations between Emperor K'ang Hsi and King John III of Poland",%
\space%
157-161.%
}
% 原书页码：321；注释 3:69
\NoteEntry{3}{69}{%
Sebes, The Jesuits and the Sino-Russian Treaty of Nerchinsk, 98.%
}
% 原书页码：321；注释 3:70
\NoteEntry{3}{70}{%
Kajdanski, “Ferdinand Verbiest's Relations with King John III of Poland”, 308。另外，莱布尼茨%
后来确实也给彼得大帝写了信，劝说俄罗斯如果能作为中西交流的桥梁，必定能获取不%
小的经济利益。Wiener,Leibniz Selection,598。莱布尼茨和闵明我神父之间的通信已部分%
翻译成中文，收录在：夏瑞春编，陈爱政等译，《德国思想家论中国》。%
}
% 原书页码：321；注释 3:71
\NoteEntry{3}{71}{%
《清代中俄档案》，1:58-59。不过荷兰帮大清带去的信在1690年才到达俄罗斯，那时中%
俄《尼布楚条约》已经签署，荷兰没帮上忙。后来大清内部以为是荷兰把信带到以后，%
才有中俄《尼布楚条约》会谈，这是错误的。见《朔方备乘》，53:1-2。%
}
% 原书页码：321；注释 3:72
\NoteEntry{3}{72}{%
《圣祖仁皇帝庭训格言》，卷1。%
}
% 原书页码：321；注释 3:73
\NoteEntry{3}{73}{%
这些法国传教士在本国的筛选过程，参看：Jami, The Emperor's New Mathematics, 105.%
}
% 原书页码：321；注释 3:74
\NoteEntry{3}{74}{%
《耶稣会士中国书简集》，1:264-265。%
}
% 原书页码：321；注释 3:75
\NoteEntry{3}{75}{%
Bouvet, The History of Cang-Hy, 52-53.%
}
% 原书页码：321；注释 3:76
\NoteEntry{3}{76}{%
Chen, "Sino-Russian Diplomatic Relations Since 1689", in The Chinese Social and Political Science%
\space%
Review，10(1926)，476-479。该刊物是民国时期北京出版的英文刊物，汇集了众多民国一%
流华人学者的英文研究。该刊所有论文在网上都可以免费查阅。该文作者陈复光从清华%
大学政治系毕业后，获得庚子赔款资助留学哈佛。学成后回国，执教于云南大学，一直%
到1949年以后。他20世纪40年代用中文写的《有清一代之中俄关系》，参考了俄语、%
拉丁语等多种西方文献，至今仍是中俄关系研究的经典。%
}
% 原书页码：321；注释 3:77
\NoteEntry{3}{77}{%
Jaskov, "The Negotiated Geography of the Treaty of Nerchinsk (1689) and the Role of the Jesuits";%
\space%
Sebes, The Jesuits and the Sino-Russian Treaty of Nerchinsk, 264.%
}
% 原书页码：321；注释 3:78
\NoteEntry{3}{78}{%
Sebes, The Jesuits and the Sino-Russian Treaty of Nerchinsk, 271. 徐日升神父的日记是用葡萄牙文%
写的，该日记的葡萄牙文和英文翻译对照，在Sebes书的最后部分。%
}
% 原书页码：321；注释 3:79
\NoteEntry{3}{79}{%
Jaskov, "The Negotiated Geography of the Treaty of Nerchinsk (1689) and the Role of the Jesuits",%
\space%
66.%
}
% 原书页码：321；注释 3:80
\NoteEntry{3}{80}{%
Sebes, The Jesuits and the Sino-Russian Treaty of Nerchinsk, 271.%
}
% 原书页码：321；注释 3:81
\NoteEntry{3}{81}{%
Sebes, The Jesuits and the Sino-Russian Treaty of Nerchinsk, 271.%
}
% 原书页码：321；注释 3:82
\NoteEntry{3}{82}{%
Jaskov, “The Negotiated Geography of the Treaty of Nerchinsk (1689) and the Role of the Jesuits”,%
\space%
71, Note 78.%
}
% 原书页码：321；注释 3:83
\NoteEntry{3}{83}{%
Jaskov, "The Negotiated Geography of the Treaty of Nerchinsk (1689) and the Role of the Jesuits”,%
\space%
72.%
}
% 原书页码：321；注释 3:84
\NoteEntry{3}{84}{%
Jaskov, “The Negotiated Geography of the Treaty of Nerchinsk (1689) and the Role of the Jesuits”%
66.%
}
% 原书页码：321；注释 3:85
\NoteEntry{3}{85}{%
Jaskov, "The Negotiated Geography of the Treaty of Nerchinsk (1689) and the Role of the Jesuits",%
\space%
70-72.%
}
% 原书页码：321；注释 3:86
\NoteEntry{3}{86}{%
Jaskov, "The Negotiated Geography of the Treaty of Nerchinsk (1689) and the Role of the Jesuits",%
\space%
73-74.%
}
% 原书页码：322；注释 3:87
\NoteEntry{3}{87}{%
Jaskov, "The Negotiated Geography of the Treary of Nerchinsk (1689) and the Role of the Jesuits",%
\space%
73-74.%
}
% 原书页码：322；注释 3:88
\NoteEntry{3}{88}{%
Chen, "Sino-Russian Diplomatic Relations Since 1689", 482-484.%
}
% 原书页码：322；注释 3:89
\NoteEntry{3}{89}{%
Jaskov, "The Negotiated Geography of the Treaty of Nerchinsk (1689) and the Role of the Jesuits”,%
\space%
64-65.%
}
% 原书页码：322；注释 3:90
\NoteEntry{3}{90}{%
Sebes, The Jesuits and the Sino-Russian Treaty of Nerchinsk, 103–104.%
}
% 原书页码：322；注释 3:91
\NoteEntry{3}{91}{%
Sebes, The Jesuits and the Sino-Russian Treaty of Nerchinsk, 109.%
}
% 原书页码：322；注释 3:92
\NoteEntry{3}{92}{%
俄罗斯代表最早是希望从传教士那里知道中国的谈判底牌。Chen,“Sino-Russian%
\space%
Diplomatic Relations Since 1689", 502.%
}
% 原书页码：322；注释 3:93
\NoteEntry{3}{93}{%
Chen, "Sino-Russian Diplomatic Relations Since 1689", 177-179.%
}
% 原书页码：322；注释 3:94
\NoteEntry{3}{94}{%
《圣祖仁皇帝庭训格言》，717:650。白晋在他1699年出版的康熙传记中，也写到康熙%
是在杨光先参劾汤若望时期开始学习西学的。Bouvet, The History of Cang-Hy, 52.%
}
% 原书页码：322；注释 3:95
\NoteEntry{3}{95}{%
昭梿，《啸亭杂录》，476。昭梿（1776—1829）是康亲王杰书的后人。%
}
% 原书页码：322；注释 3:96
\NoteEntry{3}{96}{%
《清中前期西洋天主教在华活动档案史料》，1:57。%
}
% 原书页码：322；注释 3:97
\NoteEntry{3}{97}{%
何秋涛，《朔方备乘》，1:23a。%
}
% 原书页码：322；注释 3:98
\NoteEntry{3}{98}{%
Sebes, The Jesuits and the Sino-Russian Treaty of Nerchinsk, 279.%
}
% 原书页码：322；注释 3:99
\NoteEntry{3}{99}{%
《朔方备乘》，53:1-2。%
}
% 原书页码：322；注释 3:100
\NoteEntry{3}{100}{%
Chen, "Sino-Russian Diplomatic Relations Since 1689", 507-508.%
}
% 原书页码：322；注释 3:101
\NoteEntry{3}{101}{%
陈复光，《十八世纪初叶清廷进攻准噶尔期间第一次到俄属及俄京的中国使节》；李齐芳，%
《清雍正皇帝两次遣使赴俄之谜》；Hsü,“Russia's Special Position in China during the Early%
\space%
Ch'ing Period".%
}
\NoteChapter{第四章\quad 内务府的人}
% 原书页码：322；注释 4:1
\NoteEntry{4}{1}{%
《耶稣会士中国书简集》，1:282。%
}
% 原书页码：322；注释 4:2
\NoteEntry{4}{2}{%
Witek, "Understanding the Chinese: A Comparison of Matteo Ricci and the French Jesuit%
\space%
Mathematicians Sent by Louis XIV", 90-93.%
}
% 原书页码：322–323；注释 4:3
\NoteEntry{4}{3}{%
容教诏书尚未在任何清宫档案中发现。现知最早一份写有容教诏书全文内容的文件，%
是徐日升神父一份折子的手抄副本。该抄本没有注明抄录时间，它存放在雍正时期的%
军机处档案中，很可能是在雍正初期抄录，当时朝廷在讨论禁止天主教的问题。这一%
时期，耶稣会传教士在一封信中说，他们私下重金贿赂了一位礼部文员，让他抄写了%
有关康熙容教诏书的文件。参看：《耶稣会士中国书简集》，2:324。该抄本现存于内阁%
大库资料，可网络调阅：http://archive.ihp.sinica.edu.tw/mctkm2/index.html，文档编号为%
048450。%
另外一份早期的容教诏书记录，全文见于徐日升墓碑背面的碑文。该碑文记载，石%
碑刻于1692年，是容教诏书颁布的那一年。不过也有可能是徐日升1708年去世后刻写%
的。石碑碑文在1914年做过拓本。拓本高184厘米，宽87厘米，藏于北京的中国国家%
图书馆。容教诏书全文在碑文的最后一部分。碑文的前面部分撰写了自汤若望以来耶稣%
会传教士对康熙宫廷的贡献。（因碑损，拓本前半部分已不清晰。）这块碑在1949年后%
损毁遗失。该拓片的图片已经出版，见：《北京图书馆藏中国历代石刻拓本汇编》，65:3。%
对于其他版本的容教诏书的考证，参看：Standaert, “The Edict ofTolerance: A Textual History%
\space%
and Reading", 308-358.%
}
% 原书页码：323；注释 4:4
\NoteEntry{4}{4}{%
1701年，据统计有140名欧洲传教士在中国，大概有 20万信徒。Standahert, Handbook of%
\space%
Christianity in China, 301-303, 383.%
}
% 原书页码：323；注释 4:5
\NoteEntry{4}{5}{%
胡祥雨，《“逃人法”人“顺治律”考——兼谈“逃人法”的应用》。%
}
% 原书页码：323；注释 4:6
\NoteEntry{4}{6}{%
《满文老档》，165。%
}
% 原书页码：323；注释 4:7
\NoteEntry{4}{7}{%
Spence, Ts’ao Yin and the K'ang-hsi Emperor, 10-11.%
}
% 原书页码：323；注释 4:8
\NoteEntry{4}{8}{%
Chu, Law and Society in Traditional China, 9, 128, 280-283; Sommer, Sex, Law, and Society in%
\space%
Late Imperial China, 5-8.%
}
% 原书页码：323；注释 4:9
\NoteEntry{4}{9}{%
陈国栋，《清代内务府包衣三旗人人员及其旗下组织——兼论一些有关包衣的问题》，%
325-343。%
}
% 原书页码：323；注释 4:10
\NoteEntry{4}{10}{%
Chang, "The Economic Role of the Imperial Household in the Ching Dynasty", 243-273; Feng,%
"The Emperor's Coffer: The Qing Imperial Fiscal Separation Between Privy Purse and State%
\space%
Treasury (1644-1912)", 138-168.%
}
% 原书页码：323；注释 4:11
\NoteEntry{4}{11}{%
祁美琴，《清代内务府》。%
}
% 原书页码：323；注释 4:12
\NoteEntry{4}{12}{%
Chang, "The Economic Role of the Imperial Household in the Ching Dynasty", 251.%
}
% 原书页码：323；注释 4:13
\NoteEntry{4}{13}{%
杜家骥，《清代八旗官制与行政》，171:193。%
}
% 原书页码：323；注释 4:14
\NoteEntry{4}{14}{%
Wu, "Foreword", in Torbert, Preston M., The Ch'ing Imperial Household Department, vii ; Spence,%
\space%
Ts'ao Yin and the K'ang-hsi Emperor, 134. 王钟翰在给祁美琴《清代内务府》写的序言中，也%
提出内务府和官僚机构是一种并存关系。%
}
% 原书页码：323；注释 4:15
\NoteEntry{4}{15}{%
“王者无私”出自《汉书·文帝纪》。%
}
% 原书页码：323；注释 4:16
\NoteEntry{4}{16}{%
关于明代皇帝个人生活的限制，参看《万历十五年》。%
}
% 原书页码：323；注释 4:17
\NoteEntry{4}{17}{%
傅克东，《八旗户籍制度初探》，34-35; Spence, Ts'ao Yin and the K'ang-hsi Emperor, 10-11。%
}
% 原书页码：323；注释 4:18
\NoteEntry{4}{18}{%
Stumpf, The Acta Pekinensia, or, Historical Records of the Maillard de Tournon Legation: First%
\space%
Transcribed Edition and English Annotated Translation, Volume 1. 183. Note 406。 以后注释中,%
该书（中文名《北京纪事》）简称AP。该书的第二卷（Leiden:Brill，2019）在注释中称%
为AP2。在AP书中，用“mancipia”来翻译“包衣”。当然很有可能耶稣会士在其他的%
写作中还用过别的词语来翻译“包衣”。在AP的书末，还附带了一张CD，里面有AP拉%
丁文全文PDF文档，该拉丁文文档在以后注释中称为拉丁AP。对于“mancipia”这个%
词在AP中的应用，参看拉丁AP，81，132，138，139，198，210，279，317，319，338，%
344,359, 365,366,394,398,399。%
《北京纪事》（Acta Pekinensia）是由耶稣会神父纪理安编订，由北京大部分耶稣会%
神父成员署名确认其真实性，是一部记录教皇使节到北京后各种事情的备忘录。该备忘%
录用拉丁文书写，在教皇使节离开不久后，就发回了欧洲。这份备忘录一直存放在耶稣%
会士罗马档案中。直到2015年，它的第一部分（截至1706年8月）被翻译成了英文出%
版。关于《北京纪事》的介绍，以及整个翻译项目的前前后后，参看：Rule,“The Acta%
\space%
Pekinensia Project", 17-29.%
}
% 原书页码：324；注释 4:19
\NoteEntry{4}{19}{%
AP, 198.%
}
% 原书页码：324；注释 4:20
\NoteEntry{4}{20}{%
AP, 525.%
}
% 原书页码：324；注释 4:21
\NoteEntry{4}{21}{%
《雍亲王致年羹尧书》全文见附录一。%
}
% 原书页码：324；注释 4:22
\NoteEntry{4}{22}{%
孟光祖是旗人，是皇三子的人。他到地方各省结识督抚，后来被抓。被抓后，给他定的%
罪名是骗子，是打着皇三子的名义出来行骗的。不过孟光祖是否确为骗子，还是皇三子%
一派为了摆脱干系，安插罪名来和他撇清关系，现在很难知晓。正史中记载，不只是年%
羹尧，还有好几个地方大员都和孟光祖有牵连。《圣祖仁皇帝实录》，卷279。%
}
% 原书页码：324；注释 4:23
\NoteEntry{4}{23}{%
拉丁 AP,81;AP,78.%
}
% 原书页码：324；注释 4:24
\NoteEntry{4}{24}{%
拉丁 AP, 319; AP, 529.%
}
% 原书页码：324；注释 4:25
\NoteEntry{4}{25}{%
陈垣，《释奴才》，603-607。%
}
% 原书页码：324；注释 4:26
\NoteEntry{4}{26}{%
以上选自《关于江宁织造曹家档案史料》。%
}
% 原书页码：324；注释 4:27
\NoteEntry{4}{27}{%
参看《康熙朝汉文朱批奏折汇编》。%
}
% 原书页码：324；注释 4:28
\NoteEntry{4}{28}{%
侯寿昌，《辽东佟氏族属旗籍考辨》，367。%
}
% 原书页码：324；注释 4:29
\NoteEntry{4}{29}{%
《清初西洋传教士满文档案译本》，317-318。%
}
% 原书页码：324；注释 4:30
\NoteEntry{4}{30}{%
传教士是否能算作包衣奴才？根据我与中外学者的私下交流来看，这个问题还有争议。%
争议的原因在于学界尚未在何者算作包衣奴才这个问题上达成一致。一部分学者认为，%
包衣奴才必须有严格的档案和奴籍记录支持。比如曹寅家，就在《八旗满洲氏族通谱》%
中有明确记载（感谢李文益博士为我提供资料）。目前尚未在八旗档案中发现传教士的%
记录，因而传教士不能算作包衣奴才。另外一些学者认为，在没有档案记录的情况下，%
可以通过行为方式来认定包衣奴才身份。通过行为来认定的逻辑，与今天如果没有找到%
结婚证，可以通过双方的实际行为方式来认定婚姻关系相类似。为了避免不必要的争论，%
本书认为康熙是把传教士当作包衣奴才在对待。这一说法，相当于今天我们说某某人把%
某小孩当作儿子在对待（意指双方实际行为中反映出了父子关系，但同时也承认尚无户%
口资料或者出生证明来表明二者确实是父子关系）。%
}
% 原书页码：324；注释 4:31
\NoteEntry{4}{31}{%
《清初西洋传教士满文档案译本》，312，313-316，321，324。%
}
% 原书页码：324；注释 4:32
\NoteEntry{4}{32}{%
AP2, 275.%
}
% 原书页码：324；注释 4:33
\NoteEntry{4}{33}{%
《康熙朝汉文朱批奏折汇编》，6:108。%
}
% 原书页码：324；注释 4:34
\NoteEntry{4}{34}{%
《康熙朝满文朱批奏折全译》，12。%
}
% 原书页码：324；注释 4:35
\NoteEntry{4}{35}{%
《耶稣会士中国书简集》，1:294。%
}
% 原书页码：324；注释 4:36
\NoteEntry{4}{36}{%
704年11月，康熙把总管内务府大臣找来，说准备从自已私人账户上，也就是“内库”%
里，借一万两银子给传教士修教堂。而且康熙还专门强调不收利息，借给传教士八年。%
一万两银子在当时算是一笔巨款。南怀仁做钦天监监正的时候，一年的俸禄是一百两银%
子，康熙可能也感觉钱借的比较多，还向总管内务府大臣解释了一下他为什么要借款。%
康熙说：“伊等皆为外国之人，除朕之外，又有何人照顾伊等？”见《清初西洋传教士满%
文档案译本》，283。%
}
% 原书页码：324；注释 4:37
\NoteEntry{4}{37}{%
《清初西洋传教士满文档案译本》，323。%
}
% 原书页码：324；注释 4:38
\NoteEntry{4}{38}{%
《清初西洋传教士满文档案译本》，281。%
}
% 原书页码：324；注释 4:39
\NoteEntry{4}{39}{%
《清初西洋传教士满文档案译本》，283。%
}
% 原书页码：325；注释 4:40
\NoteEntry{4}{40}{%
AP, 112.%
}
% 原书页码：325；注释 4:41
\NoteEntry{4}{41}{%
AP, 347.%
}
% 原书页码：325；注释 4:42
\NoteEntry{4}{42}{%
各种称谓见于《关于江宁织造曹家档案史料》。%
}
% 原书页码：325；注释 4:43
\NoteEntry{4}{43}{%
Hsia, "Some Observations on the Observations: the Decline of the French Jesuit Scientific Mission%
\space%
in China", 309.%
}
% 原书页码：325；注释 4:44
\NoteEntry{4}{44}{%
Witek, "Understanding the Chinese", 72.%
}
% 原书页码：325；注释 4:45
\NoteEntry{4}{45}{%
《耶稣会士中国书简集》，1:146-148。%
}
% 原书页码：325；注释 4:46
\NoteEntry{4}{46}{%
《耶稣会士中国书简集》，1:147。%
}
% 原书页码：325；注释 4:47
\NoteEntry{4}{47}{%
《耶稣会士中国书简集》，1:290。%
}
% 原书页码：325；注释 4:48
\NoteEntry{4}{48}{%
和太监的友善关系，见《耶稣会士中国书简集》，1:294；Ripa, Memoirs of Father Ripa, 47-%
48.%
}
% 原书页码：325；注释 4:49
\NoteEntry{4}{49}{%
《耶稣会士中国书简集》，1:313-315。%
}
% 原书页码：325；注释 4:50
\NoteEntry{4}{50}{%
《耶稣会士中国书简集》，1:315。%
}
% 原书页码：325；注释 4:51
\NoteEntry{4}{51}{%
《耶稣会士中国书简集》，1:324-327。%
}
% 原书页码：325；注释 4:52
\NoteEntry{4}{52}{%
Jami, The Emperor's New Mathematics, 102-119.%
}
% 原书页码：325；注释 4:53
\NoteEntry{4}{53}{%
《清初西洋传教士满文档案译本》，326。%
}
% 原书页码：325；注释 4:54
\NoteEntry{4}{54}{%
《耶稣会士中国书简集》，2:106。%
}
% 原书页码：325；注释 4:55
\NoteEntry{4}{55}{%
《耶稣会士中国书简集》，2:90。%
}
% 原书页码：325；注释 4:56
\NoteEntry{4}{56}{%
Hanson, and Pomata, "Travels of a Chinese Pulse Treatise: The Latin and French Translations of the%
\space%
Tuzhu maijue bianzhen圖註脈訣辨真(1650s-1730s)”, 24.%
}
% 原书页码：325；注释 4:57
\NoteEntry{4}{57}{%
Szczesniak, "John Floyer and Chinese Medicine", 133-138.%
}
% 原书页码：325；注释 4:58
\NoteEntry{4}{58}{%
Puente-Ballesteros, "Jesuit Medicine in the Kangxi Court (1662-1722): Imperial Networks and%
\space%
Patronage", 144-145.%
}
% 原书页码：325；注释 4:59
\NoteEntry{4}{59}{%
《清初西洋传教士满文档案译本》，297。%
}
% 原书页码：325；注释 4:60
\NoteEntry{4}{60}{%
《耶稣会士中国书简集》，4:159。%
}
% 原书页码：325；注释 4:61
\NoteEntry{4}{61}{%
《耶稣会士中国书简集》，2:305-306。%
}
% 原书页码：325；注释 4:62
\NoteEntry{4}{62}{%
《耶稣会士中国书简集》，2:308-309。%
}
% 原书页码：325；注释 4:63
\NoteEntry{4}{63}{%
《耶稣会士中国书简集》，2:310-311。%
}
% 原书页码：325；注释 4:64
\NoteEntry{4}{64}{%
Menegon, "Beijing as a Missionary Translation Center in the Eighteenth Century", 37–74.%
}
% 原书页码：325；注释 4:65
\NoteEntry{4}{65}{%
Curtis, “A Plan of the Emperor's Glassworks”, 81-90.%
}
% 原书页码：325；注释 4:66
\NoteEntry{4}{66}{%
Puente-Ballesteros, “Jesuit Medicine in the Kangxi Court (1662-1722): Imperial Networks and%
\space%
Patronage", 143.%
}
% 原书页码：325；注释 4:67
\NoteEntry{4}{67}{%
《清初西洋传教士满文档案译本》，327。%
}
% 原书页码：325；注释 4:68
\NoteEntry{4}{68}{%
该药的功能和为什么安多的病上要使用，参看：Puente-Ballesteros,“Antoine Thomas, SI As a%
\space%
Patient' of the Kangxi Emperor (r. 1662-1722): A Case Study on the Appropriation of Theriac at%
\space%
the Imperial Court", 213-250.%
}
% 原书页码：325；注释 4:69
\NoteEntry{4}{69}{%
刘世珣，《底野迦的再现：康熙年间宫廷西药德里鸦噶初探》。%
}
% 原书页码：326；注释 4:70
\NoteEntry{4}{70}{%
《关于江宁织造曹家档案史料》，98-99。%
}
% 原书页码：326；注释 4:71
\NoteEntry{4}{71}{%
《耶稣会士中国书简集》，1:282。%
}
\NoteChapter{第五章\quad 教皇使团和康熙派出的洋钦差}
% 原书页码：326；注释 5:1
\NoteEntry{5}{1}{%
Dai, The Sichuan Frontier and Tibet, 44.%
}
% 原书页码：326；注释 5:2
\NoteEntry{5}{2}{%
Antonucci, "Nores on the Qing-Zunghar War (1690-1697): The Jesuit Sources", 371-387.%
}
% 原书页码：326；注释 5:3
\NoteEntry{5}{3}{%
Tuttle, "Tibetan Buddhism at Wutai Shan in the Qing", 171-173.%
}
% 原书页码：326；注释 5:4
\NoteEntry{5}{4}{%
Farquhar, "Emperor as Bodhisattva in the Governance of The Ch'ing Empire", 23.%
}
% 原书页码：326；注释 5:5
\NoteEntry{5}{5}{%
Dai, The Sichuan Frontier and Tibet, 47.%
}
% 原书页码：326；注释 5:6
\NoteEntry{5}{6}{%
以前关于这个使团的研究，参看：Rouleau,“Maillard de Tournon”, 269; Collani,“Kangxi's%
\space%
Mandate of Heaven and Papal Authority”, 185-192; Menegon, "A Clash of Court Cultures: Papal%
\space%
Envoys in Early Eighteenth Century Beijing", 141-142.%
}
% 原书页码：326；注释 5:7
\NoteEntry{5}{7}{%
本章主要是通过对比研读当时传教士留下的记录和清宫档案来完成的。清宫档案主要来%
源：《康熙朝满文朱批奏折全译》与《清初西洋传教士满文档案译本》。%
}
% 原书页码：326；注释 5:8
\NoteEntry{5}{8}{%
《清初西洋传教士满文档案译本》，284-285；AP，92；AP2，36.%
}
% 原书页码：326；注释 5:9
\NoteEntry{5}{9}{%
该圣旨全文见：Rosso, Apostolic Legations, 112-116.%
}
% 原书页码：326；注释 5:10
\NoteEntry{5}{10}{%
《清初西洋传教士满文档案译本》,284。%
}
% 原书页码：326；注释 5:11
\NoteEntry{5}{11}{%
传教士内部确实对要不要接待多罗有过分歧。北京传教士开会达成了一致，参看：Rosso，%
Apostolic Legations, 156. 后来在多罗到京前一个多月，北京传教士在1705年10月31 日又%
开过一次会，并在这次会上对于如何会见多罗达成了一致。参看：Collani,“Kilian Srumpf%
\space%
and his Acta Pekinensia: Life, Context and Purpose of a Manuscript", 65.%
}
% 原书页码：326；注释 5:12
\NoteEntry{5}{12}{%
《清初西洋传教士满文档案译本》，284。%
}
% 原书页码：326；注释 5:13
\NoteEntry{5}{13}{%
《清初西洋传教士满文档案译本》，284。%
}
% 原书页码：326；注释 5:14
\NoteEntry{5}{14}{%
《清初西洋传教士满文档案译本》，285；AP，15。%
}
% 原书页码：326；注释 5:15
\NoteEntry{5}{15}{%
《清初西洋传教士满文档案译本》，286。%
}
% 原书页码：326；注释 5:16
\NoteEntry{5}{16}{%
《清初西洋传教士满文档案译本》，287。%
}
% 原书页码：326；注释 5:17
\NoteEntry{5}{17}{%
《清初西洋传教士满文档案译本》，289。%
}
% 原书页码：326；注释 5:18
\NoteEntry{5}{18}{%
《清初西洋传教士满文档案译本》，289。%
}
% 原书页码：326；注释 5:19
\NoteEntry{5}{19}{%
《清初西洋传教士满文档案译本》，290。%
}
% 原书页码：326；注释 5:20
\NoteEntry{5}{20}{%
AP2, 37.%
}
% 原书页码：326；注释 5:21
\NoteEntry{5}{21}{%
AP, 81.%
}
% 原书页码：326；注释 5:22
\NoteEntry{5}{22}{%
AP, 83.%
}
% 原书页码：326；注释 5:23
\NoteEntry{5}{23}{%
AP, 84.%
}
% 原书页码：326；注释 5:24
\NoteEntry{5}{24}{%
AP, 85-86.%
}
% 原书页码：327；注释 5:25
\NoteEntry{5}{25}{%
AP, 63.%
}
% 原书页码：327；注释 5:26
\NoteEntry{5}{26}{%
AP，88.建立与大清的联系是多罗来华的一项使命。他到广州后，在那里买了一处房子，%
作为罗马与大清的专用联络处。参看：Menegon,“Interlopers at the Fringes of Empirc”,%
\space%
36-37.%
}
% 原书页码：327；注释 5:27
\NoteEntry{5}{27}{%
AP, 88.%
}
% 原书页码：327；注释 5:28
\NoteEntry{5}{28}{%
AP, 89-90.%
}
% 原书页码：327；注释 5:29
\NoteEntry{5}{29}{%
AP, 90-92.%
}
% 原书页码：327；注释 5:30
\NoteEntry{5}{30}{%
传信部和多罗来华的关系，参看：Menegon, “Interlopers at the Fringes of Empire: The Procura-%
tors of the Propaganda Fide Papal Congregation in Canton and Macao, 1700-1823", 29-33.%
}
% 原书页码：327；注释 5:31
\NoteEntry{5}{31}{%
AP, 101-102.%
}
% 原书页码：327；注释 5:32
\NoteEntry{5}{32}{%
AP, 111,116.%
}
% 原书页码：327；注释 5:33
\NoteEntry{5}{33}{%
AP, 112.%
}
% 原书页码：327；注释 5:34
\NoteEntry{5}{34}{%
AP, 378.%
}
% 原书页码：327；注释 5:35
\NoteEntry{5}{35}{%
AP, 398.%
}
% 原书页码：327；注释 5:36
\NoteEntry{5}{36}{%
《康熙朝满文朱批奏折全译》，448-449。%
}
% 原书页码：327；注释 5:37
\NoteEntry{5}{37}{%
AP, 122.%
}
% 原书页码：327；注释 5:38
\NoteEntry{5}{38}{%
AP, 157-170.%
}
% 原书页码：327；注释 5:39
\NoteEntry{5}{39}{%
三个石碑具体内容参看：Jérôme Tobar, “Inscriptions juives de K'ai-Fong-Fou”.%
}
% 原书页码：327；注释 5:40
\NoteEntry{5}{40}{%
Chen, "Buddhist Monk or Jewish Rabbi?", 142.%
}
% 原书页码：327；注释 5:41
\NoteEntry{5}{41}{%
Pollak, The Torah Scroll of the Chinese Jews, 90.%
}
% 原书页码：327；注释 5:42
\NoteEntry{5}{42}{%
英国图书馆对此有专文介绍：https://blogs.bl.uk/asian-and-african/2020/06/an-eighth-century-%
judaco-persian-letter-from-dandan-uiliq.html.%
}
% 原书页码：327；注释 5:43
\NoteEntry{5}{43}{%
开封犹太人，参看： Leslie, The Survival of the Chinese Jews; Xu, The Jews of Kaifeng, China: History,%
\space%
Culture and Religion.%
}
% 原书页码：327；注释 5:44
\NoteEntry{5}{44}{%
《耶稣会士中国书简集》，2:11-20。%
}
% 原书页码：327；注释 5:45
\NoteEntry{5}{45}{%
AP, 156.%
}
% 原书页码：327；注释 5:46
\NoteEntry{5}{46}{%
关于中印两国礼仪之争的争论，最新研究总结参看：Županov,The Rites Controversies in the%
\space%
Early Modern World.%
}
% 原书页码：327；注释 5:47
\NoteEntry{5}{47}{%
AP, 141.%
}
% 原书页码：327；注释 5:48
\NoteEntry{5}{48}{%
AP, 202.%
}
% 原书页码：327；注释 5:49
\NoteEntry{5}{49}{%
前后过程，参看：Minamiki, The Chinese Rites Controversy from Its Beginning to Modern Times, 25-%
42.%
}
% 原书页码：327；注释 5:50
\NoteEntry{5}{50}{%
AP, 141.%
}
% 原书页码：327；注释 5:51
\NoteEntry{5}{51}{%
AP, 12-13.%
}
% 原书页码：327；注释 5:52
\NoteEntry{5}{52}{%
《清初西洋传教士满文档案译本》，No.33。%
}
% 原书页码：327–328；注释 5:53
\NoteEntry{5}{53}{%
对于《北京纪事》，有一位拉丁语读者写过一篇读后感。他在其中写到，他是一个懂%
拉丁文，但对中国历史一无所知的人。所以他的读后感更接近当时欧洲人读过该报告%
后的感觉。Hughes, “The Acta's Four Portraits”, 89-138.%
}
% 原书页码：328；注释 5:54
\NoteEntry{5}{54}{%
AP, 141.%
}
% 原书页码：328；注释 5:55
\NoteEntry{5}{55}{%
AP, 181-197.%
}
% 原书页码：328；注释 5:56
\NoteEntry{5}{56}{%
《康熙朝满文朱批奏折全译》，1577-1578。%
}
% 原书页码：328；注释 5:57
\NoteEntry{5}{57}{%
《清初西洋传教士满文档案译本》，294。%
}
% 原书页码：328；注释 5:58
\NoteEntry{5}{58}{%
AP，285. 关于北京当时的水资源管理，参看：Lei, “One City with Two Waters: Drinking%
\space%
Water in Beijing, 1644–1900", 483-484.%
}
% 原书页码：328；注释 5:59
\NoteEntry{5}{59}{%
AP, 407-409.%
}
% 原书页码：328；注释 5:60
\NoteEntry{5}{60}{%
AP, 408.%
}
% 原书页码：328；注释 5:61
\NoteEntry{5}{61}{%
AP, 411.%
}
% 原书页码：328；注释 5:62
\NoteEntry{5}{62}{%
AP, 412.%
}
% 原书页码：328；注释 5:63
\NoteEntry{5}{63}{%
AP, 412.%
}
% 原书页码：328；注释 5:64
\NoteEntry{5}{64}{%
AP, 414.%
}
% 原书页码：328；注释 5:65
\NoteEntry{5}{65}{%
AP, 417.%
}
% 原书页码：328；注释 5:66
\NoteEntry{5}{66}{%
《康熙朝满文朱批奏折全译》，418；《清初西洋传教士满文档案译本》，297。%
}
% 原书页码：328；注释 5:67
\NoteEntry{5}{67}{%
《康熙朝满文朱批奏折全译》，419-421。%
}
% 原书页码：328；注释 5:68
\NoteEntry{5}{68}{%
《康熙朝满文朱批奏折全译》，421。%
}
% 原书页码：328；注释 5:69
\NoteEntry{5}{69}{%
关于颜珰在礼仪之争中的角色，参看：Collani,“Charles Maigrot's Role in the Chinese Rites%
\space%
Controversy", 149-183.%
}
% 原书页码：328；注释 5:70
\NoteEntry{5}{70}{%
《康熙朝满文朱批奏折全译》，422；《清初西洋传教士满文档案译本》，300。%
}
% 原书页码：328；注释 5:71
\NoteEntry{5}{71}{%
AP, 690-691.%
}
% 原书页码：328；注释 5:72
\NoteEntry{5}{72}{%
AP, 691.%
}
% 原书页码：328；注释 5:73
\NoteEntry{5}{73}{%
《康熙朝满文朱批奏折全译》，435。%
}
% 原书页码：328；注释 5:74
\NoteEntry{5}{74}{%
AP, 583-586.%
}
% 原书页码：328；注释 5:75
\NoteEntry{5}{75}{%
《康熙朝满文朱批奏折全译》，435。%
}
% 原书页码：328；注释 5:76
\NoteEntry{5}{76}{%
AP, 4.%
}
% 原书页码：328；注释 5:77
\NoteEntry{5}{77}{%
《清初西洋传教士满文档案译本》，308。%
}
% 原书页码：328；注释 5:78
\NoteEntry{5}{78}{%
AP，5，注释 5.%
}
% 原书页码：328；注释 5:79
\NoteEntry{5}{79}{%
《康熙朝满文朱批奏折全译》，462。%
}
% 原书页码：328；注释 5:80
\NoteEntry{5}{80}{%
AP, 5.%
}
% 原书页码：328；注释 5:81
\NoteEntry{5}{81}{%
《清初西洋传教士满文档案译本》，308；AP2，124-125.%
}
% 原书页码：328；注释 5:82
\NoteEntry{5}{82}{%
AP2, 269.%
}
% 原书页码：328；注释 5:83
\NoteEntry{5}{83}{%
《康熙朝满文朱批奏折全译》，435。%
}
% 原书页码：328；注释 5:84
\NoteEntry{5}{84}{%
《正教奉褒》，364；类似文字还见于：《清初西洋传教士满文档案译本》，311。%
}
% 原书页码：328；注释 5:85
\NoteEntry{5}{85}{%
《清初西洋传教士满文档案译本》，314。%
}
% 原书页码：328；注释 5:86
\NoteEntry{5}{86}{%
AP2, 275.%
}
% 原书页码：329；注释 5:87
\NoteEntry{5}{87}{%
关于多罗谕令，参看：Malatesta, “A Fatal Clash of Wllls”.%
}
% 原书页码：329；注释 5:88
\NoteEntry{5}{88}{%
《中国礼仪之争：西文文献一百篇（1645—1941）》，文献8。译文有少许改动。%
}
% 原书页码：329；注释 5:89
\NoteEntry{5}{89}{%
《康熙与罗马使节关系文书》，文件14。%
}
% 原书页码：329；注释 5:90
\NoteEntry{5}{90}{%
AP2, 377-380, 395-397.%
}
% 原书页码：329；注释 5:91
\NoteEntry{5}{91}{%
Fairbank; Tëng, "On The Ch’ing Tributary System", 135-246.%
}
% 原书页码：329；注释 5:92
\NoteEntry{5}{92}{%
AP2, 107-108.%
}
% 原书页码：329；注释 5:93
\NoteEntry{5}{93}{%
《康熙朝满文朱批奏折全译》，421。%
}
\NoteChapter{第六章\quad 彼得大帝使团和教皇新使团}
% 原书页码：329；注释 6:1
\NoteEntry{6}{1}{%
Malatesta, "A Fatal Clash of Wills", 219-225; AP2, 304-308.%
}
% 原书页码：329；注释 6:2
\NoteEntry{6}{2}{%
AP2, 321.%
}
% 原书页码：329；注释 6:3
\NoteEntry{6}{3}{%
徐日升神父写给南方神父的信，已经翻译成英文，参看：Malatesta,“A Fatal Clash of Wills”，%
242-245; AP2, 331-333.%
}
% 原书页码：329；注释 6:4
\NoteEntry{6}{4}{%
Malatesta, "A Fatal Clash of Wills", 240; AP2, 333-334.%
}
% 原书页码：329；注释 6:5
\NoteEntry{6}{5}{%
张诚写给多罗的信，参看：Malatesta, “A Fatal Clash of Wills”, 234-239; AP2, 268-269. 对于张%
诚葬礼的分析，参看：Collani, “From the Earthly Court to the Heavenly Court", 128-130.%
}
% 原书页码：329；注释 6:6
\NoteEntry{6}{6}{%
Rosso, Apostolic Legations, 156.%
}
% 原书页码：329；注释 6:7
\NoteEntry{6}{7}{%
AP, 71.%
}
% 原书页码：329；注释 6:8
\NoteEntry{6}{8}{%
AP, 70.%
}
% 原书页码：329；注释 6:9
\NoteEntry{6}{9}{%
张诚写给多罗的信全文见：Malatesta, “A Fatal Clash of Wills”, 234-239.%
}
% 原书页码：329；注释 6:10
\NoteEntry{6}{10}{%
从康熙1706年派出使节出使欧洲，到1720年第二次接待教皇特使，再到1721年特使离%
开，康熙在欧洲问题上花了很多心思。以前有人说康熙在与欧洲打交道的漫长过程中，%
丧失了与欧洲交流的兴趣。再加上雍正皇帝即位后一年多，于1724年1月禁止了天主教%
在华传播。这段康熙与欧洲的交往就被认为是天主教在华传播失败的导火索。前人总结%
天主教在清初被禁的原因时，也把禁教与康熙联系到一起。雍正的禁教被说成是康熙晚%
年天主教政策的延续。这是当下学界流行了几十年的解释。关于雍正禁教源于康熙的论%
述，始于 1940 年代的 Rowbotham, Missionary and Mandarin, 176；六十年后出版的基督教传%
华史研究手册，也坚持该论。见 Standaert, Handbook of Christianity in China, 498-499, 519。%
国内学界也把海外研究中这个说法当成定论。比如新近的中文研究中，陈青松《关于康%
熙朝一条查禁天主教史料的重新认识》文中第一段就明言：“由于礼仪之争’中罗马教%
廷的蛮横态度，康熙帝下令禁教，这已经成为学界的共识。”由于此论在中外学界流行了%
八十年，这里就不一一列举。最新近的两篇：Hsia,“Imperial China and Jesuit Mission”，44;%
\space%
Collani, "The Jesuit Rites Controversy".%
}
% 原书页码：329；注释 6:11
\NoteEntry{6}{11}{%
Rouleau, "Maillard de Tournon", 265-266, Note 3.%
}
% 原书页码：329；注释 6:12
\NoteEntry{6}{12}{%
Witek, "Eliminating Misunderstandings", 211-247.%
}
% 原书页码：329；注释 6:13
\NoteEntry{6}{13}{%
《康熙朝满文朱批奏折全译》，544。%
}
% 原书页码：330；注释 6:14
\NoteEntry{6}{14}{%
Rosso, Apostolic Legations, 199.%
}
% 原书页码：330；注释 6:15
\NoteEntry{6}{15}{%
《康熙朝满文朱批奏折全译》，586。%
}
% 原书页码：330；注释 6:16
\NoteEntry{6}{16}{%
《康熙朝满文朱批奏折全译》，594。%
}
% 原书页码：330；注释 6:17
\NoteEntry{6}{17}{%
《康熙朝汉文朱批奏折汇编》，2:334,364,370-371,380-382。%
}
% 原书页码：330；注释 6:18
\NoteEntry{6}{18}{%
《康熙朝汉文朱批奏折汇编》，2:385-597,410,439-441。%
}
% 原书页码：330；注释 6:19
\NoteEntry{6}{19}{%
《康熙朝满文朱批奏折全译》，633，634，650。%
}
% 原书页码：330；注释 6:20
\NoteEntry{6}{20}{%
《康熙朝汉文朱批奏折汇编》，2:761。%
}
% 原书页码：330；注释 6:21
\NoteEntry{6}{21}{%
《康熙朝汉文朱批奏折汇编》，3:5。%
}
% 原书页码：330；注释 6:22
\NoteEntry{6}{22}{%
《康熙朝汉文朱批奏折汇编》，3:6-11。%
}
% 原书页码：330；注释 6:23
\NoteEntry{6}{23}{%
《康熙朝满文朱批奏折全译》，735，741。%
}
% 原书页码：330；注释 6:24
\NoteEntry{6}{24}{%
《康熙朝满文朱批奏折全译》，804。%
}
% 原书页码：330；注释 6:25
\NoteEntry{6}{25}{%
《康熙朝满文朱批奏折全译》，804-805。%
}
% 原书页码：330；注释 6:26
\NoteEntry{6}{26}{%
《康熙朝满文朱批奏折全译》，818。%
}
% 原书页码：330；注释 6:27
\NoteEntry{6}{27}{%
内阁大库档案，档案编号401001362。%
}
% 原书页码：330；注释 6:28
\NoteEntry{6}{28}{%
《康熙朝汉文朱批奏折汇编》，8:198。%
}
% 原书页码：330；注释 6:29
\NoteEntry{6}{29}{%
《康熙朝汉文朱批奏折汇编》，6:439-442,7:356-357,422,451-452。%
}
% 原书页码：330；注释 6:30
\NoteEntry{6}{30}{%
《康熙朝汉文朱批奏折汇编》，6:439。%
}
% 原书页码：330；注释 6:31
\NoteEntry{6}{31}{%
教皇这一次用了他们天主教内部最高级别的命令手段来强制执行禁约。相关分析参看：%
Minamiki, The Chinese Rites Controversy, 58-62.%
}
% 原书页码：330；注释 6:32
\NoteEntry{6}{32}{%
Rosso, Apostolic Legations, 192.%
}
% 原书页码：330；注释 6:33
\NoteEntry{6}{33}{%
Dai, The Sichuan Frontier and Tibet, 48.%
}
% 原书页码：330；注释 6:34
\NoteEntry{6}{34}{%
Rosso, Apostolic Legations, 199.%
}
% 原书页码：330；注释 6:35
\NoteEntry{6}{35}{%
信的全文已经有中文翻译，参看：《中国礼仪之争：西文文献一百篇（1645—1941）》，文%
献17。%
}
% 原书页码：330；注释 6:36
\NoteEntry{6}{36}{%
《康熙与罗马使节关系文书》，文件13。%
}
% 原书页码：330；注释 6:37
\NoteEntry{6}{37}{%
Bell, Travels from St. Petersburg, in Russia, to diverse parts of Asia, 2: 37.%
}
% 原书页码：330；注释 6:38
\NoteEntry{6}{38}{%
Spence, "Claims and Counter-claims: The Kangxi Emperor and the Europeans (1661-1722)", 22.%
}
% 原书页码：330；注释 6:39
\NoteEntry{6}{39}{%
Dawson, The Chinese Chameleon, 56.%
}
% 原书页码：330；注释 6:40
\NoteEntry{6}{40}{%
Allsop, "Teodorico Pedrini: The Music and Letters of an 18th-century Missionary in China", 57-58.%
}
% 原书页码：330；注释 6:41
\NoteEntry{6}{41}{%
Malatesta, "A Fatal Clash of Wills", 230-231; AP2, 293.%
}
% 原书页码：330；注释 6:42
\NoteEntry{6}{42}{%
AP, 4.%
}
% 原书页码：330；注释 6:43
\NoteEntry{6}{43}{%
Michele Fatica 在罗马一个档案馆中发现这份密折的抄件。抄件的全文图片收录在：Fatica,%
"Matteo Ripa Journal  (1716-1720) with New Documents for the History of the Kangxi Era",%
\space%
406-417.%
}
% 原书页码：330；注释 6:44
\NoteEntry{6}{44}{%
Paul Rule教授在一封电子邮件中告诉我此人是多罗使团中的人物，原名是 Giuseppe%
\space%
Ignazio Cordero。%
}
% 原书页码：331；注释 6:45
\NoteEntry{6}{45}{%
《康熙朝满文朱批奏折全译》，823-824。这封罗若德转来的信，康熙收到后，找传教士%
翻译成了中文。关于这封信的一些背景，参看：韩琦，《姗姗来迟的“西洋消息”一%
1709年教皇致康熙信到达宫廷始末》，1-14。%
}
% 原书页码：331；注释 6:46
\NoteEntry{6}{46}{%
《清初西洋传教士满文档案译本》，331；关于刘应对中国礼仪的看法，见：AP，310-311;%
\space%
AP2, 45-46.%
}
% 原书页码：331；注释 6:47
\NoteEntry{6}{47}{%
Tarsetti, “Mariner, musician and missionary, and true priest always: Teodorico Pedrini's llfe in%
\space%
Xitang", 130.%
}
% 原书页码：331；注释 6:48
\NoteEntry{6}{48}{%
《康熙朝满文朱批奏折全译》，1158,1160。%
}
% 原书页码：331；注释 6:49
\NoteEntry{6}{49}{%
Cams, Companions in Geography.%
}
% 原书页码：331；注释 6:50
\NoteEntry{6}{50}{%
《清初西洋传教士满文档案译本》，326。%
}
% 原书页码：331；注释 6:51
\NoteEntry{6}{51}{%
《清中前期西洋天主教在华活动档案史料》，1:23.%
}
% 原书页码：331；注释 6:52
\NoteEntry{6}{52}{%
Jin, "The Mission to Congratulate the Kangxi Emperor on his Seventieth Birthday: Father João%
\space%
Mourão and His Plan to Bring an Ambassador of the King of Portugal to Beijing", 428.%
}
% 原书页码：331；注释 6:53
\NoteEntry{6}{53}{%
《清代中俄关系档案史料选编》，400。%
}
% 原书页码：331；注释 6:54
\NoteEntry{6}{54}{%
Mancall, Russia and China, 216-217.%
}
% 原书页码：331；注释 6:55
\NoteEntry{6}{55}{%
《清代中俄关系档案史料选编》，408。%
}
% 原书页码：331；注释 6:56
\NoteEntry{6}{56}{%
Bell, Travels from St. Petersburg, in Russia, to diverse parts of Asia, 1: 356.%
}
% 原书页码：331；注释 6:57
\NoteEntry{6}{57}{%
Bell, Travels from St. Petersburg, in Russia, to diverse parts of Asia, 2: 3.%
}
% 原书页码：331；注释 6:58
\NoteEntry{6}{58}{%
Bell, Travels from St. Petersburg, in Russia, to diverse parts of Asia, 2: 5.%
}
% 原书页码：331；注释 6:59
\NoteEntry{6}{59}{%
Bell, Travels from St. Petersburg, in Russia, to diverse parts of Asia, 2: 7.%
}
% 原书页码：331；注释 6:60
\NoteEntry{6}{60}{%
整个接待的详情见： Bell, Travels from St. Petersburg, in Russia, to diverse parts of Asia, 2: 7-11.%
}
% 原书页码：331；注释 6:61
\NoteEntry{6}{61}{%
Bell, Travels from St. Petersburg, in Russia, to diverse parts of Asia, 2: 18.%
}
% 原书页码：331；注释 6:62
\NoteEntry{6}{62}{%
Bell, Travels from St. Petersburg, in Russia, to diverse parts of Asia, 2: 32–34.%
}
% 原书页码：331；注释 6:63
\NoteEntry{6}{63}{%
Bell, Travels from St. Petersburg, in Russia, to diverse parts of Asia, 2: 19–20.%
}
% 原书页码：331；注释 6:64
\NoteEntry{6}{64}{%
Bell, Travels from St. Petersburg, in Russia, to diverse parts of Asia, 2: 44.%
}
% 原书页码：331；注释 6:65
\NoteEntry{6}{65}{%
Bell, Travels from St. Petersburg, in Russia, to diverse parts of Asia, 2: 46–47.%
}
% 原书页码：331；注释 6:66
\NoteEntry{6}{66}{%
Bell, Travels from St. Petersburg, in Russia, to diverse parts of Asia, 2: 54, 73.%
}
% 原书页码：331；注释 6:67
\NoteEntry{6}{67}{%
Benedict, Golden-Silk Smoke: A History of Tobacco in China, 1550–2010, 7.%
}
% 原书页码：331；注释 6:68
\NoteEntry{6}{68}{%
Bell, Travels from St. Petersburg, in Russia, to diverse parts of Asia, 2: 102.%
}
% 原书页码：331；注释 6:69
\NoteEntry{6}{69}{%
Bell, Travels from St. Petersburg, in Russia, to diverse parts of Asia, 2: 50.%
}
% 原书页码：331；注释 6:70
\NoteEntry{6}{70}{%
Bell, Travels from St. Petersburg, in Russia, to diverse parts of Asia, 2: 71.%
}
% 原书页码：331；注释 6:71
\NoteEntry{6}{71}{%
《清代中俄关系档案史料选编》，407.%
}
% 原书页码：331；注释 6:72
\NoteEntry{6}{72}{%
《清代中俄关系档案史料选编》，407.%
}
% 原书页码：331；注释 6:73
\NoteEntry{6}{73}{%
Bell, Travels from St. Petersburg, in Russia, to diverse parts of Asia, 2: 85–86.%
}
% 原书页码：331；注释 6:74
\NoteEntry{6}{74}{%
《康熙与罗马使节关系文书》，文件12。%
}
% 原书页码：331；注释 6:75
\NoteEntry{6}{75}{%
《康熙与罗马使节关系文书》，文件13。%
}
% 原书页码：332；注释 6:76
\NoteEntry{6}{76}{%
《康熙与罗马使节关系文书》，文件13。%
}
% 原书页码：332；注释 6:77
\NoteEntry{6}{77}{%
《康熙与罗马使节关系文书》，文件13。%
}
% 原书页码：332；注释 6:78
\NoteEntry{6}{78}{%
《康熙与罗马使节关系文书》，文件13。%
}
% 原书页码：332；注释 6:79
\NoteEntry{6}{79}{%
《康熙与罗马使节关系文书》，文件13。%
}
% 原书页码：332；注释 6:80
\NoteEntry{6}{80}{%
《康熙与罗马使节关系文书》，文件13。%
}
% 原书页码：332；注释 6:81
\NoteEntry{6}{81}{%
关于康熙先让传教士和嘉乐谈，最后又决定亲自和嘉乐谈，这个转变约翰·贝尔有记录：%
Bell, Travels from St. Petersburg, in Russia, to diverse parts of Asia, 2: 47.%
}
% 原书页码：332；注释 6:82
\NoteEntry{6}{82}{%
《康熙与罗马使节关系文书》，文件13。%
}
% 原书页码：332；注释 6:83
\NoteEntry{6}{83}{%
Witek, Controversial Ideas in China and in Europe, 6.%
}
% 原书页码：332；注释 6:84
\NoteEntry{6}{84}{%
Minamiki, The Chinese Rites Controversy, 64-65.%
}
% 原书页码：332；注释 6:85
\NoteEntry{6}{85}{%
《中国礼仪之争：西文文献一百篇（1645—1941）》，文献24。%
}
% 原书页码：332；注释 6:86
\NoteEntry{6}{86}{%
《康熙与罗马使节关系文书》，文件13。%
}
% 原书页码：332；注释 6:87
\NoteEntry{6}{87}{%
这一天是康熙派人记录和嘉乐谈判的最后一天。这份记录用汉语写成，用今天的%
话说，相当于会议纪要。尽管后来嘉乐在北京一直待到过完春节，还参加了康熙的元%
宵节聚会，但这段话是整个记录的最后一段。也就是说，在康熙眼里，整个谈判在这%
里就结束了。%
}
% 原书页码：332；注释 6:88
\NoteEntry{6}{88}{%
《康熙与罗马使节关系文书》，文件13。%
}
% 原书页码：332；注释 6:89
\NoteEntry{6}{89}{%
Rosso, Apostolic Legations, 383-387.%
}
% 原书页码：332；注释 6:90
\NoteEntry{6}{90}{%
Minamiki, The Chinese Rites Controversy, 72.%
}
% 原书页码：332；注释 6:91
\NoteEntry{6}{91}{%
《康熙朝汉文朱批奏折汇编》，8:905-906。%
}
% 原书页码：332；注释 6:92
\NoteEntry{6}{92}{%
Gray, Walks in the City of Canton, 633-635.%
}
% 原书页码：332；注释 6:93
\NoteEntry{6}{93}{%
Rosso, Apostolic Legations, 209.%
}
% 原书页码：332；注释 6:94
\NoteEntry{6}{94}{%
《康熙朝汉文朱批奏折汇编》，8:912。%
}
% 原书页码：332；注释 6:95
\NoteEntry{6}{95}{%
Rosso, Apostolic Legations, 211.%
}
% 原书页码：332；注释 6:96
\NoteEntry{6}{96}{%
Bell, Travels from St. Petersburg, in Russia, to diverse parts of Asia, 2: 69–70.%
}
% 原书页码：332；注释 6:97
\NoteEntry{6}{97}{%
比如乾隆年间，官方编订介绍世界各国的《皇清职贡图》在写到意大利时，就提到了教皇，%
书中称其为“教化王”。但提到和大清的关系时，该书说教皇是雍正年间开始派使团到%
大清的，也就是不知道康熙年间的两个使团。《皇清职贡图》，1:33a。%
}
\NoteChapter{第七章\quad 雍正争位与禁教}
% 原书页码：332；注释 7:1
\NoteEntry{7}{1}{%
《雍正起居注》，2:1175-1176。%
}
% 原书页码：332；注释 7:2
\NoteEntry{7}{2}{%
Rosso, Apostolic Legations, 216-219.%
}
% 原书页码：332；注释 7:3
\NoteEntry{7}{3}{%
Ripa, Memoirs of Father Ripa, 124.%
}
% 原书页码：332；注释 7:4
\NoteEntry{7}{4}{%
AP, 415-416.%
}
% 原书页码：332；注释 7:5
\NoteEntry{7}{5}{%
Jami, The Emperor's New Mathematics, 290-291.%
}
% 原书页码：333；注释 7:6
\NoteEntry{7}{6}{%
《康熙朝满文奏折全译》，1158。%
}
% 原书页码：333；注释 7:7
\NoteEntry{7}{7}{%
Huang, Autocracy at Work, 71–75.%
}
% 原书页码：333；注释 7:8
\NoteEntry{7}{8}{%
《允禩允禟案：穆经远口供》，1:1b。%
}
% 原书页码：333；注释 7:9
\NoteEntry{7}{9}{%
Bell, Travels from St. Petersburg in Russia, to diverse parts of Asia, 2: 124.%
}
% 原书页码：333；注释 7:10
\NoteEntry{7}{10}{%
Spence, Treason by the Book, 246.%
}
% 原书页码：333；注释 7:11
\NoteEntry{7}{11}{%
孟森，《清世宗人承大统考实》，519-572。%
}
% 原书页码：333；注释 7:12
\NoteEntry{7}{12}{%
《允禩允糖案：穆经远口供》。%
}
% 原书页码：333；注释 7:13
\NoteEntry{7}{13}{%
全文刊载于 Rosso, Apostolic Legations, 407-419.%
}
% 原书页码：333；注释 7:14
\NoteEntry{7}{14}{%
《耶稣会士中国书简集》，4:66-74。%
}
% 原书页码：333；注释 7:15
\NoteEntry{7}{15}{%
王钟翰，《清世宗夺嫡考实》，147-193。%
}
% 原书页码：333；注释 7:16
\NoteEntry{7}{16}{%
满保在雍正元年七月二十九日用满语写了一封密折向雍正汇报过这件事了。雍正在批示%
中夸满保汇报得好，让他用“缮本”奏到朝廷的渠道中去。《雍正朝满文朱批奏折全译》，%
1:257-258。%
}
% 原书页码：333；注释 7:17
\NoteEntry{7}{17}{%
《耶稣会士中国书简集》，2:314-315。%
}
% 原书页码：333；注释 7:18
\NoteEntry{7}{18}{%
Menegon, Ancestors, Virgins, and Friars, 116–124.%
}
% 原书页码：333；注释 7:19
\NoteEntry{7}{19}{%
《耶稣会士中国书简集》，2:322。%
}
% 原书页码：333；注释 7:20
\NoteEntry{7}{20}{%
《耶稣会士中国书简集》，2:323。%
}
% 原书页码：333；注释 7:21
\NoteEntry{7}{21}{%
《耶稣会士中国书简集》，2:323。%
}
% 原书页码：333；注释 7:22
\NoteEntry{7}{22}{%
《耶稣会士中国书简集》，2:324。%
}
% 原书页码：333；注释 7:23
\NoteEntry{7}{23}{%
《耶稣会士中国书简集》，2:324-325。%
}
% 原书页码：333；注释 7:24
\NoteEntry{7}{24}{%
《耶稣会士中国书简集》，2:327。%
}
% 原书页码：333；注释 7:25
\NoteEntry{7}{25}{%
《耶稣会士中国书简集》，2:334-335。%
}
% 原书页码：333；注释 7:26
\NoteEntry{7}{26}{%
《耶稣会士中国书简集》，2:334-335。%
}
% 原书页码：333；注释 7:27
\NoteEntry{7}{27}{%
《清中前期西洋天主教在华活动档案史料》，1:58-59。%
}
% 原书页码：333；注释 7:28
\NoteEntry{7}{28}{%
《耶稣会士中国书简集》，2:337,340。%
}
% 原书页码：333；注释 7:29
\NoteEntry{7}{29}{%
《雍正朝满文朱批奏折全译》，1:258。%
}
% 原书页码：333；注释 7:30
\NoteEntry{7}{30}{%
《耶稣会士中国书简集》，2:324-325。%
}
% 原书页码：333；注释 7:31
\NoteEntry{7}{31}{%
《清中前期西洋天主教在华活动档案史料》，1:59。%
}
% 原书页码：333；注释 7:32
\NoteEntry{7}{32}{%
《宫中档雍正朝奏折》，3:393。%
}
% 原书页码：333；注释 7:33
\NoteEntry{7}{33}{%
《耶稣会士中国书简集》，2:322。%
}
% 原书页码：333；注释 7:34
\NoteEntry{7}{34}{%
《耶稣会士中国书简集》，2:335。%
}
% 原书页码：333；注释 7:35
\NoteEntry{7}{35}{%
欧美学者的论述：Farquhar, “Emperor as Bodhisattva”, 32; Rockhill, “The Dalai Lamas of Lhasa%
\space%
and Their Relations with the Manchu Emperors of China, 1644–1908”; Petech, China and Tibet%
\space%
in the Early 18th Century; Hevia, "Lamas, Emperors, and Rituals”, 243-278; Schwieger, The Dalai%
\space%
Lama and the Emperor of China, 112-145.%
}
% 原书页码：333；注释 7:36
\NoteEntry{7}{36}{%
孟森，《清代史》，257;冯尔康，《雍正传》，505-517。%
}
% 原书页码：334；注释 7:37
\NoteEntry{7}{37}{%
杭世骏，《武林理安寺志》，243-244。%
}
% 原书页码：334；注释 7:38
\NoteEntry{7}{38}{%
《御选语录》，854-855。%
}
% 原书页码：334；注释 7:39
\NoteEntry{7}{39}{%
《御选语录》，855。%
}
% 原书页码：334；注释 7:40
\NoteEntry{7}{40}{%
对该活佛在北京的活动，参看：Wang, “The Qing Court's Tibet Connection”, 125-163.%
}
% 原书页码：334；注释 7:41
\NoteEntry{7}{41}{%
《清世宗关于佛学之谕旨》，4:8b-9a。%
}
% 原书页码：334；注释 7:42
\NoteEntry{7}{42}{%
虚云，《上海玉佛寺禅七开示》。%
}
% 原书页码：334；注释 7:43
\NoteEntry{7}{43}{%
印光，《拣魔辨异录石印序》，卷3。%
}
% 原书页码：334；注释 7:44
\NoteEntry{7}{44}{%
Menegon, "Yongzheng's Conundrum. The Emperor on Christianity, Religions, and Heterodoxy", 330.%
}
% 原书页码：334；注释 7:45
\NoteEntry{7}{45}{%
《康熙与罗马使节关系文书》，39。%
}
% 原书页码：334；注释 7:46
\NoteEntry{7}{46}{%
Rosso, Apostolic Legations, 214.%
}
% 原书页码：334；注释 7:47
\NoteEntry{7}{47}{%
对于苏努一家的研究，参看：陈垣，《雍正间奉天主教之宗室》；方豪，《清代旗人之信奉%
天主教与遭禁》。%
}
% 原书页码：334；注释 7:48
\NoteEntry{7}{48}{%
王冕森，《允禟西洋密码字档案初探》。%
}
% 原书页码：334；注释 7:49
\NoteEntry{7}{49}{%
《耶稣会士中国书简集》，3:26-27。%
}
% 原书页码：334；注释 7:50
\NoteEntry{7}{50}{%
Tarsetti, "Mariner, musician and missionary, and true priest always: Teodorico Pedrini's life in%
\space%
Xitang", 147.%
}
% 原书页码：334；注释 7:51
\NoteEntry{7}{51}{%
《允裸允案：穆经远口供》。%
}
% 原书页码：334；注释 7:52
\NoteEntry{7}{52}{%
曹树基估计清代初年有一亿六千万人口，清代年人口增长率为千分之七，据此推算，所%
以雍正年间约有两亿人口。曹树基，《中国人口史（第五卷清时期）》，51。%
}
% 原书页码：334；注释 7:53
\NoteEntry{7}{53}{%
根据1668年的一个统计数据。周叔迦，《清代佛教史料辑稿》，88。%
}
% 原书页码：334；注释 7:54
\NoteEntry{7}{54}{%
杜文凯，《清代西人见闻录》，145。%
}
% 原书页码：334；注释 7:55
\NoteEntry{7}{55}{%
《世宗宪皇帝实录》，卷3。%
}
% 原书页码：334；注释 7:56
\NoteEntry{7}{56}{%
《雍正起居注》，1:466。%
}
% 原书页码：334；注释 7:57
\NoteEntry{7}{57}{%
Perdue, China Marches West, 245.%
}
\NoteChapter{跋}
% 原书页码：334；注释 epilogue:1
\NoteEntry{epilogue}{1}{%
Jami, The Emperor's New Mathematics, 288-294.%
}
% 原书页码：334；注释 epilogue:2
\NoteEntry{epilogue}{2}{%
罗见今，《明安图是卡塔兰数的首创者》，239-245; Larcombe,“The 18th Century Chinese%
\space%
Discovery of the Catalan Numbers", 5-6。%
}
% 原书页码：334；注释 epilogue:3
\NoteEntry{epilogue}{3}{%
比如，席泽宗院士一篇影响大的文章就专门引用到了这段话。参看：席泽宗，《论康熙科%
学政策的失误》，21。%
}
% 原书页码：334；注释 epilogue:4
\NoteEntry{epilogue}{4}{%
Katz, and Parshall, Taming the Unknown: A History of Algebra from Antiquity to the Early Twentieth%
\space%
Century, 277-279.%
}
% 原书页码：335；注释 epilogue:5
\NoteEntry{epilogue}{5}{%
Jami, The Emperor's New Mathematics, 300-301.%
}
% 原书页码：335；注释 epilogue:6
\NoteEntry{epilogue}{6}{%
《康熙朝满文奏折全译》，1158,1160。%
}
% 原书页码：335；注释 epilogue:7
\NoteEntry{epilogue}{7}{%
潘亦宁，《韦达方程解法在康熙时代的传播》，365。%
}
% 原书页码：335；注释 epilogue:8
\NoteEntry{epilogue}{8}{%
Jami, The Emperor's New Mathematics, 299.%
}
% 原书页码：335；注释 epilogue:9
\NoteEntry{epilogue}{9}{%
比如: Hao, Reception of Du Fu (712-770) and His Poetry in Imperial China.%
}
% 原书页码：335；注释 epilogue:10
\NoteEntry{epilogue}{10}{%
Du Jarric, Akbar and the Jesuits: an account of the Jesuit missions to the court of Akbar.%
}
% 原书页码：335；注释 epilogue:11
\NoteEntry{epilogue}{11}{%
Pomplun, Jesuit on the Roof of the World: Ippolito Desideri's Mission to Tibet.%
}
% 原书页码：335；注释 epilogue:12
\NoteEntry{epilogue}{12}{%
早期传教士往亚洲传教，在决定传教地点时，分析思路确实类似今天对市场的分析。顺%
治时代，去过中国西藏、尼泊尔等地的传教士，就认为西藏难度大，传教机会小，但是%
他们认为尼泊尔有机会，因为那里的统治者不是和尚。Givon,“Connecting Eurasia: Jesuit%
\space%
Experimentation with Overland Mobility Between China and Europe, 1656-1664”, 664.%
}
% 原书页码：335；注释 epilogue:13
\NoteEntry{epilogue}{13}{%
Musillo, "Reconciling Two Careers".%
}
% 原书页码：335；注释 epilogue:14
\NoteEntry{epilogue}{14}{%
1760年代，传教士遇到财产分配等问题，是找内务府官员来断他们的是非，而不是找官%
府。Hsia, “Jesuit Survival and Restoration in China”, 251-256.%
}
% 原书页码：335；注释 epilogue:15
\NoteEntry{epilogue}{15}{%
关于金尼阁的自杀，以及后来耶稣会内部的调查，参看：Logan and Brockey,“Nicolas%
\space%
Trigault, S. J.: A Portrait by Peter Paul Rubens", 161.%
}
\endgroup
